% Options for packages loaded elsewhere
\PassOptionsToPackage{unicode}{hyperref}
\PassOptionsToPackage{hyphens}{url}
%
\documentclass[titlepage,twoside,openright,11pt]{article}
\usepackage{amsmath,amssymb}
\usepackage{iftex}
\ifPDFTeX
  \usepackage[T1]{fontenc}
  \usepackage[utf8]{inputenc}
  \usepackage{textcomp} % provide euro and other symbols
\else % if luatex or xetex
  \usepackage{unicode-math} % this also loads fontspec
  \defaultfontfeatures{Scale=MatchLowercase}
  \defaultfontfeatures[\rmfamily]{Ligatures=TeX,Scale=1}
\fi
\usepackage{lmodern}
\ifPDFTeX\else
  % xetex/luatex font selection
\fi
% Use upquote if available, for straight quotes in verbatim environments
\IfFileExists{upquote.sty}{\usepackage{upquote}}{}
\IfFileExists{microtype.sty}{% use microtype if available
  \usepackage[]{microtype}
  \UseMicrotypeSet[protrusion]{basicmath} % disable protrusion for tt fonts
}{}
\makeatletter
\@ifundefined{KOMAClassName}{% if non-KOMA class
  \IfFileExists{parskip.sty}{%
    \usepackage{parskip}
  }{% else
    \setlength{\parindent}{0pt}
    \setlength{\parskip}{6pt plus 2pt minus 1pt}}
}{% if KOMA class
  \KOMAoptions{parskip=half}}
\makeatother
\usepackage{xcolor}
\usepackage[letterpaper,inner=1.25in,outer=1in,top=1in,bottom=1in]{geometry}
\setlength{\emergencystretch}{3em} % prevent overfull lines
\providecommand{\tightlist}{%
  \setlength{\itemsep}{0pt}\setlength{\parskip}{0pt}}
\setcounter{secnumdepth}{2}
\setcounter{tocdepth}{2}
% Make each top-level \section start on a new right-hand page, so that
% sections open like chapters in a printed book.
\usepackage{titlesec}
\newcommand{\sectionbreak}{\cleardoublepage}
\ifLuaTeX
  \usepackage{selnolig}  % disable illegal ligatures
\fi
\usepackage{bookmark}
\IfFileExists{xurl.sty}{\usepackage{xurl}}{} % add URL line breaks if available
\urlstyle{same}
\hypersetup{
  pdftitle={Scone User's Guide},
  pdfauthor={Scott E. Fahlman},
  colorlinks=true,
  linkcolor=blue!50!black,
  urlcolor=blue!60!black,
  citecolor=blue!50!black,
  pdfcreator={LaTeX via pandoc}}

\title{Scone User's Guide}
\author{Scott E.\ Fahlman\\[6pt]
  Language Technologies Institute \&\\
  Computer Science Department\\
  Carnegie Mellon University\\[6pt]
  \texttt{sef@cs.cmu.edu}}
\date{Version 1.0.0 \\[4pt] August 2014}

\begin{document}
\maketitle

This manual is copyright © 2003-2014 by Scott E. Fahlman.

Carnegie Mellon University holds the copyright to the core Scone
software. The Scone software is made available to the public under the
Apache 2.0 open source license. A copy of this license is distributed
with the software. The license can also be found at
\url{http://www.apache.org/licenses/LICENSE-2.0}.

Unless required by applicable law or agreed to in writing, software
distributed under the License is distributed on an "AS IS" BASIS,
WITHOUT WARRANTIES OR CONDITIONS OF ANY KIND, either express or implied.
See the License for the specific language governing permissions and
limitations under the License.

Acknowledgments:

Development of Scone since February 2013 has been supported in part by
the U.S. Office of Naval Research under award number N000141310224.

Development of Scone from January 2012 through August 2013 was supported
in part by the Intelligence Advanced Research Projects Activity (IARPA)
via Department of Defense US Army Research Laboratory contract number
W911NF-12-C-0020.

Development of Scone from 2003 through 2008 was supported in part by the
Defense Advanced Research Projects Agency (DARPA) under contract numbers
NBCHD030010 and FA8750-07-D-0185.

Additional support for Scone development has been provided by research
grants from Cisco Systems Inc. and from Google Inc.

The U.S. Government is authorized to reproduce and distribute reprints
for Governmental purposes notwithstanding any copyright annotation
thereon.

Disclaimer: The views and conclusions contained herein are those of the
authors and should not be interpreted as necessarily representing the
official policies or endorsements, either expressed or implied, of
IARPA, DoD/ARL, ONR, DARPA, the U.S. Government, or any of our other
sponsors.

\cleardoublepage
\tableofcontents
\cleardoublepage

\section{Overview}\label{overview}

\href{http://www.cs.cmu.edu/~sef/scone}{Scone} is a \emph{knowledge-base
system} (KBS) intended for use as a component in a wide range of
software applications. The Scone system comprises a search/inference
engine and a ``core'' \emph{knowledge base} (KB) of common knowledge.
Application-specific KBs and interfaces to other software applications
can be built on top of this common base.

This manual is intended as a \emph{programmer-level} guide to the
operation and functions of the Scone software. It is not intended to be
a tutorial on knowledge representation concepts used in Scone. A
book-length tutorial document is needed, and it is currently under
development. Some of the ideas in Scone are discussed in essay-length
posts on Scott Fahlman's ``Knowledge Nuggets'' blog:
\url{http://www.cs.cmu.edu/~nuggets/}. Readers interested in the
intellectual origins of Scone should also consult Scott E.\ Fahlman's
book \emph{NETL: A System for Representing and Using Real-World
Knowledge} (MIT Press, 1979), which introduced the marker-passing,
parallel-inheritance style of knowledge representation that Scone
continues and extends. Scone is currently under active
development. We will include an updated version of this manual with each
release, along with other release notes as appropriate.

The greatest problem for users of knowledge-base systems has been the
difficulty of adding new knowledge to the system and making that
knowledge fully effective. This gives rise to spotty coverage of
common-sense domains. Scone attempts to ease this burden by relatively
clean design, expressiveness of the representation system, and by
insulating the user from most efficiency concerns. In addition, the
Scone group has developed a number of tools for converting knowledge
from other structured or semi-structured formats (existing ontologies
such as SUMO and WordNet, tables, databases, and information mined from
the Internet) to Scone format. Extraction of Scone knowledge from
natural-language text, such as Wikipedia or newswire stories, is a
long-term goal and is the subject of active research in our group.

Scone is written in Common Lisp.\footnote{It would be possible to
  re-implement the Scone engine in some other language such as Java,
  perhaps with certain performance-critical code in C for efficiency.
  However, we believe that Common Lisp is by far the best language for
  the \emph{development} of Scone, and (with careful coding) the Lisp
  version is nearly as efficient as it would be in C. The downside is
  that experienced Lisp programmers are becoming an endangered species.}
The primary vehicle for Scone development is
\href{http://www.sbcl.org/}{Steel Bank Common Lisp} (SBCL), which is a
derivative of \href{http://www.cons.org/cmucl/}{CMU Common Lisp}
(CMUCL). SBCL is a free, high-performance implementation of Common Lisp
running on Linux and some other operating systems. The SBCL system makes
use of the 64-bit address space on Intel and AMD processors, making it
possible to run very large knowledge bases on those machines.

Scone also has been run successfully on CMU Common Lisp and on several
Windows-based Common Lisp systems, including CLISP, LispWorks, and
Allegro Common Lisp. However, as of now, we do not actively test or
maintain Scone in those environments.

Scone can be run as a server process on an Intel/Linux system,
communicating with other software applications via character-stream I/O
using the TCP/IP socket mechanism. The other software applications may
be written in any computer language and may be running on the same
machine or a different machine on the Internet.

Scone is designed for efficiency of search and inference on conventional
high-end workstations. However, the marker-passing inference algorithms
used in Scone were originally designed with parallelism in mind. Scone
is therefore well suited for future implementation on a massively
parallel machine or perhaps a computing grid.

\subsubsection{Note on Common Lisp
Syntax}\label{note-on-common-lisp-syntax}

This manual assumes that the reader has some general familiarity with
Common Lisp\footnote{For a good introduction to Common Lisp, we suggest
  \href{http://www.gigamonkeys.com/book/}{Practical Common Lisp} by
  Peter Seibel, available for free online or in hardcopy form at
  booksellers. For excellent examples of how to use Lisp in AI
  programming, including some knowledge representation examples, see
  \href{http://norvig.com/paip.html}{Paradigms of AI Programming} by
  Peter Norvig. The definitive online reference for Common Lisp is the
  \href{http://www.lispworks.com/documentation/HyperSpec/Front/}{Common
  Lisp Hyperspec}, produced by Kent M. Pitman for Lispworks Ltd., based
  on the ANSI X3J13 language standard.} syntax. As a quick reminder, if
this document indicates that a function's syntax looks like this

\textbf{some-function (a b \&optional c)} \emph{{[}FUNCTION{]}}

then \textbf{some-function} is the name of the function, arguments
\textbf{a} and \textbf{b} are required, while argument \textbf{c} is
optional, with some default value. Unless otherwise specified, the
default default value is \textbf{nil} -- Lisp\textquotesingle s
representation for both the empty list and "false".

If a function\textquotesingle s syntax specification looks like this

\textbf{some-function (a b \&key :cat :dog)} \emph{{[}FUNCTION{]}}

then \textbf{a} and \textbf{b} are required arguments, while
\textbf{:dog} and \textbf{:cat} are optional ``keyword'' arguments,
which the caller can pass in any order or not at all. Some syntactically
legal calls to this function would include the following

\textbf{(some-function 27 \textquotesingle my-arg)\\
(some-function 27 \textquotesingle my-arg :cat "Felix" :dog "Fido")\\
(some-function 27 \textquotesingle my-arg :dog "Fido" :cat "Felix")}

In Common Lisp, \textbf{\textquotesingle foo} represents a quoted
constant value (in this case the symbol \textbf{foo}), and
\textbf{"foo"} represents a string object. Symbols prefixed with the
colon character are called \emph{keywords}, and do not have to be
quoted; these are used for named tokens or enumerations and for labeling
the keyword/value pairs in a function call.

In Common Lisp, functions may return more than one value via the
\textbf{values} form. This is useful when a function generates more than
one useful value -- for example, to return both a quotient and a
remainder from a division function. If the function is called normally,
return values beyond the first one are discarded. However, if the
function is called under \textbf{multiple-value-bind} or similar
functions, the additional return values are captured.

The curly-brace \textbf{\{something\}} syntax that Scone uses for
element internal names is not standard Common Lisp syntax; it is a
Scone-specific extension created using Lisp\textquotesingle s character
macro facility.

\section{Getting Started}\label{getting-started}

Scone is not a stand-alone executable. It is loaded as source code into
a running Common Lisp image, and all interaction happens at the Lisp
top level (the ``REPL''). This chapter walks through starting the
system, loading a knowledge base, and shutting down. The chapters that
follow describe the available functions in detail.

Scone requires a Common Lisp implementation. It is primarily developed
and tested with Steel Bank Common Lisp (SBCL), and has also been run on
CMUCL, CLISP, LispWorks, and Allegro Common Lisp. There is no build
step.

\subsection{Starting Scone}\label{starting-scone}

The recommended way to start Scone is via \textbf{scone-loader.lisp}.
From the top level of the Scone source tree, start your Lisp and
evaluate:

\begin{verbatim}
(load "scone-loader.lisp")
(scone)
\end{verbatim}

\textbf{(load "scone-loader.lisp")} reads the loader file and prints a
reminder of how to invoke it. \textbf{(scone)} then loads
\textbf{engine.lisp}, sets \textbf{*default-kb-pathname*} to point at
this repository's \textbf{kb/} directory, and loads the bootstrap KB --
the structures the engine itself refers to (\{thing\}, \{set\}, and so
on). This bootstrap load is mandatory; most engine functions will not
work without it.

At this point Scone is running and you can start creating elements,
running queries, or loading additional KB files as described below.

\paragraph{Alternative: loading engine.lisp directly.} If you prefer
not to use the loader, you can perform the same steps by hand:

\begin{verbatim}
(load "engine.lisp")

(setf *default-kb-pathname*
      (merge-pathnames "kb/anonymous.lisp" (truename ".")))

(load-kb "bootstrap")
\end{verbatim}

\textbf{*default-kb-pathname*} tells \textbf{load-kb} where to look for
KB files. Set it once to any file inside your \textbf{kb/} directory
(the filename part is overridden on each \textbf{load-kb} call) and
subsequent loads can refer to KB files by short name. The specific file
\textbf{anonymous.lisp} does not need to exist; only the directory part
of the path matters.

\subsection{Loading a Knowledge Base}\label{loading-a-knowledge-base}

Scone knowledge-base files are plain Common Lisp source files with the
extension \textbf{.lisp}. Load them with \textbf{load-kb}, passing a
path relative to \textbf{*default-kb-pathname*} (the ``\textbf{.lisp}''
extension is supplied automatically). The knowledge bases shipped with
Scone fall into three groups, and for each group the call looks like:

\begin{itemize}
\item
  \textbf{core} -- the generic common-sense core ontology (upper
  ontology, units, space, time, simple physics, biology, information,
  social model, person model, episodic model, and geopolitics). Most
  applications will want this:
  \begin{verbatim}
(load-kb "core")
  \end{verbatim}
\item
  \textbf{core-components/*} -- the individual files that make up the
  core. Load one of these singly only if you want a subset of the core
  instead of the full \textbf{core} bundle above:
  \begin{verbatim}
(load-kb "core-components/time-model")
  \end{verbatim}
\item
  \textbf{lexical-components/*} -- WordNet-derived concept and name
  libraries. These are very large and take a while to load; pull them
  in only if you need broad lexical coverage:
  \begin{verbatim}
(load-kb "lexical-components/wordnet-concepts")
  \end{verbatim}
\end{itemize}

Your own application KB is just another Lisp file that uses Scone's
element-creation forms:

\begin{verbatim}
(load-kb "my-application-kb")
\end{verbatim}

KB files may call \textbf{load-kb} on other KB files, so a single
top-level load can pull in an entire dependency tree. Scone records
every file it has loaded in the variable \textbf{*loaded-files*}.

\textbf{load-kb} does not guard against loading the same KB twice. If
you have already loaded \textbf{core}, do not follow it with, for
example, \textbf{(load-kb "core-components/time-model")} -- that file
was already loaded transitively, and reloading it will signal an
``iname already in use'' error as Scone tries to redefine elements
that already exist.

\subsubsection*{Dependencies between KBs}

The shipped knowledge bases are \emph{additive}, not mutually
exclusive, but they are layered -- each one builds on the ones before
it. The overall stack is:

\begin{itemize}
\item
  \textbf{bootstrap} -- mandatory, always loaded first.
\item
  \textbf{core-components/*} -- loaded in this fixed order by
  \textbf{core.lisp}:
  \textbf{upper} $\rightarrow$ \textbf{units} $\rightarrow$
  \textbf{space-model} $\rightarrow$ \textbf{time-model} $\rightarrow$
  \textbf{simple-physics} $\rightarrow$ \textbf{simple-biology}
  $\rightarrow$ \textbf{simple-information} $\rightarrow$
  \textbf{simple-social-model} $\rightarrow$
  \textbf{simple-person-model} $\rightarrow$
  \textbf{simple-episodic-model} $\rightarrow$
  \textbf{simple-geopolitics}.
  Each file freely references concepts defined in the earlier ones.
\item
  \textbf{lexical-components/*} -- optional add-ons that layer on top
  of a loaded core. They are independent of one another, so you can
  load only the ones you need (for example, English names without the
  Spanish names).
\item
  \textbf{your application KBs} -- any further KB files you write,
  loaded on top of whatever base you have established.
\end{itemize}

Nothing in a core-component file calls \textbf{load-kb} on the files
it conceptually depends on, so the contract ``load predecessors first''
is enforced only by convention. The practical consequences:

\begin{itemize}
\item
  \textbf{(load-kb "core")} pulls in all eleven core-components in the
  right order. This is the common case.
\item
  You may load a \emph{prefix} of the chain on its own -- for example,
  just \textbf{upper} + \textbf{units} + \textbf{space-model} -- if you
  want a lean KB.
\item
  You may \emph{not} skip earlier components and load a later one on
  its own. \textbf{simple-person-model}, for instance, references
  concepts from \textbf{simple-biology}; loading it without the
  predecessors will either fail or silently create \{undefined thing\}
  stubs for the missing concepts.
\end{itemize}

The full set of load-related functions, variables, and error-checking
controls is documented in the next chapter.

\subsection{Saving and Exiting}\label{saving-and-exiting}

Scone has no formal ``shutdown'' procedure: when you exit your Lisp,
the in-memory knowledge base is discarded. To preserve work done in a
running session, \emph{checkpoint} the KB to a text file first:

\begin{verbatim}
(checkpoint-kb "my-session")
\end{verbatim}

This writes a ``\textbf{.lisp}'' file (under
\textbf{*default-kb-pathname*}) containing every element added since
the bootstrap file was loaded. In a fresh Lisp you can reload the
session with:

\begin{verbatim}
(load "engine.lisp")
(setf *default-kb-pathname*
      (merge-pathnames "kb/anonymous.lisp" (truename ".")))
(load-kb "bootstrap")
(load-kb "my-session")
\end{verbatim}

Checkpointing, incremental logging, and the related controls are
covered in detail in the next chapter.

To exit your Lisp, use its normal quit command:

\begin{itemize}
\tightlist
\item
  \textbf{SBCL}: \textbf{(exit)} or \textbf{(sb-ext:exit)}
\item
  \textbf{CMUCL, CCL, CLISP}: \textbf{(quit)}
\item
  \textbf{LispWorks}: \textbf{(lw:quit)} or close the IDE window
\end{itemize}

If you expect to resume the same Scone session repeatedly, many Lisps
can save a pre-loaded memory image that restarts far faster than
re-loading the engine and KB files from source. In SBCL, for example,
after loading engine and KBs:

\begin{verbatim}
(sb-ext:save-lisp-and-die "scone.core")
\end{verbatim}

Then start a subsequent session with \textbf{sbcl
-\/-core scone.core}. Consult your Lisp implementation's manual for
the equivalent facility.

\section{Loading, Saving, and Checkpointing KB
Files}\label{loading-saving-and-checkpointing-kb-files}

The knowledge-base (or ``KB'') files in Scone are just files of Common
Lisp expressions and comments, following the syntax of Common Lisp.
Users can easily modify these files offline, using the text editor of
their choice. These files have the extension ``\textbf{.lisp}''. Since
these text-format KB files are just Common Lisp source files, any Common
Lisp expressions you like may be included in these files, mixed with the
Scone-specific expressions that create new elements.

\subsubsection{Functions and Variables}\label{functions-and-variables}

\paragraph{Loading KB Files}\label{loading-kb-files}

\textbf{load-kb} (filename \&key :verbose) \emph{{[}FUNCTION{]}}\\
Loads the named file, with missing parts of the pathname filled in from
the default in \textbf{*default-kb-pathname*}. This is just a wrapper on
Lisp\textquotesingle s built-in \textbf{load} function, but
\textbf{load-kb} does some internal bookkeeping so that the most recent
loads can be rolled back if necessary, and so that the system knows it
is in non-interactive mode during the load.

Often the KB file will effectively be a script that loads a list of
other KB files. If the \textbf{:verbose} argument is \textbf{t}, the
result of loading and executing each form in the file will be printed.

Example of usage: \textbf{(load-kb ``project1/geopolitics'' :verbose t)}

\textbf{*default-kb-pathname*} \emph{{[}VARIABLE{]}}\\
A string. The default location for text-format KB files. This is
initialized by the \textbf{scone} function, which starts some version of
Scone.

\textbf{*loading-kb-file*} \emph{{[}VARIABLE{]}}\\
If reading a KB file, this will be \textbf{t}. Certain interactive
dialogues may be suppressed, etc.

\textbf{*loaded-files*} \emph{{[}VARIABLE{]}}\\
A list of files loaded in the current KB. The head of the list is the
file loaded most recently. Files go on this list when they have been
loaded completely.

\textbf{*currently-loading-files*} \emph{{[}VARIABLE{]}}\\
A list of the files currently being loaded. One file may load another,
and so on, so there many be a number of files on this list. The first
file on the list is the most recently started (i.e. the load that is
most deeply nested), while the last file is generally the one for which
the user called \textbf{load-kb}.

\textbf{*verbose-loading*} \emph{{[}VARIABLE{]}}\\
If set, loading files will produce printout showing progress. Normally
set by the \textbf{:verbose} keyword argument in \textbf{load-kb}.

\textbf{*disambiguate-policy*} \emph{{[}VARIABLE{]}}\\
This global variable controls what happens when an ambiguous element
name is encountered during input. \textbf{:error} says to signal an
error. \textbf{:ask} says to query the user interactively.

\paragraph{Variables Controlling Inference and Checking During KB
Modification}\label{variables-controlling-inference-and-checking-during-kb-modification}

\textbf{*reloading-kb-file*} \emph{{[}VARIABLE{]}}\\
Normally when loading forms into the KB, Scone will do extra work to see
if the new knowledge is legal and consistent, and whether Scone must
create derived elements implied by the new knowledge. If we are
re-loading a KB file that was dumped from Scone, we set
\textbf{*reloading-kb-file*} to \textbf{t}, indicating that all this
checking has already been done and does not need to be repeated. This
can greatly speed up the re-loading. Note that if other KB files have
changed significantly since the dump occurred, it is best to load,
rather than re-load, the dump -- that is, we want to do the checking
after all.

\textbf{*no-kb-error-checking*} \emph{{[}VARIABLE{]}}\\
If \textbf{t}, suppress most error checking when creating KB elements.
Use this only when loading a file that you\textquotesingle ve loaded
successfully before. This speeds up loading, but not as much as
*reloading-kb-file*. The latter also skips looking for derived
information, assuming that this is all present in the dump file.

\textbf{*check-defined-types*} \emph{{[}VARIABLE{]}}\\
If \textbf{t} (the default), check whether newly defined elements fit
the definition of some defined type, and, if so, make it an instance or
subtype of that defined-type. If \textbf{nil}, don\textquotesingle t
bother.

\textbf{*create-undefined-elements*} \emph{{[}VARIABLE{]}}\\
If \textbf{t} (the default), check whether newly defined elements fit
the definition of some defined type. If \textbf{nil},
don\textquotesingle t bother.

\textbf{*deduce-owner-type-from-role*} \emph{{[}VARIABLE{]}}\\
If \textbf{nil}, a call to x-is-the-y-of-z (or a variant) signals an
error if Z does not have or inherit a Y-role. If t (the default), try to
create an is-a link from Z to the owner of role-node Y. If successful, Y
is now inherited and we can go ahead and give it value X. If not, due to
a type mismatch or some other error, return NIL and do not modify the
KB.

\textbf{*comment-on-element-creation*} \emph{{[}VARIABLE{]}}\\
If \textbf{t} (the default), whenever Scone creates inferred KB
structures (nodes or link) not specifically requested by the user, it
prints a comment about this via the commentary stream. This can become
verbose if, for example, we are loading a file that defines types in
some random order, rather than top-down, but it also catches lots of
errors and typos.

\textbf{*comment-on-defined-types*} \emph{{[}VARIABLE{]}}\\
If \textbf{t}, whenever Scone detects that some element satisfies the
definition of a defined type and adds the element to that type, it
prints a comment about this via the commentary stream. This is useful
for debugging, but can be verbose when loading files that are known to
be correct, so the default is \textbf{nil}.

\paragraph{Deferred Connections}\label{deferred-connections}

Note: The deferred-connection mechanism has been largely superseded by
the newer mechanism that creates an \textbf{\{unknown thing\}} element
whenever the user makes reference to an element that has not yet been
explicitly defined. The deferred-connection mechanism is now used only
for bootstrapping of the root elements of the knowledge base, and we
hope to remove it altogether in the near future.

Normally we build a Scone knowledge base in orderly layers. A new
element will be connected to other elements that have already been
created. But it some situations it is necessary to create forward
references or circular references. In such cases, we use the ``deferred
connections'' mechanism. This facility is enabled by placing
\textbf{(setq *defer-unknown-connections* t)} at the start of a file
that contains forward references. Normally this variable is
\textbf{nil}, and references to unknown elements signal an error.

When \textbf{*defer-unknown-connections*} is \textbf{t} and an
element-creation routine is told to connect a wire to some element,
specified by its iname, and if no such element exists at present in the
KB, we push an entry onto the \textbf{*deferred-connections*} list so
that we can try to create the connection later, after additional
elements have been created or loaded. This allows us to handle
situations with forward or circular references, even if the
circularities span multiple file boundaries.

\textbf{*defer-unknown-connections*} \emph{{[}VARIABLE{]}}\\
If \textbf{t}, this enables the deferred-connection mechanism. This
allows for references to elements that have not yet been defined, but
that is not always a good thing: many such references are errors. So
this is \textbf{nil} by default, and we only activate it for files that
actually need this.

\textbf{*deferred-connections*} \emph{{[}VARIABLE{]}}\\
This is the list of connections and element-names that we should try to
set up later, after the named elements have been defined.

\textbf{process-deferred-connections} nil \emph{{[}FUNCTION{]}}\\
The \textbf{process-deferred-connections} function says that we should
scan the \textbf{*deferred-connections*} list now to see if any of the
deferred connections can now be created. Any connections successfully
processed are removed from the \textbf{*deferred-connections*} list; the
rest remain on the list. This function is called after loading each new
KB file, but users may occasionally want to call it directly.

\textbf{defer-connection} (tag source target) \emph{{[}FUNCTION{]}}\\
Record a deferred connection: the wire named \textbf{tag} (\textbf{:a},
\textbf{:b}, \textbf{:c}, \textbf{:parent}, \textbf{:context}, or
\textbf{:split}) on element \textbf{source} should later be connected
to \textbf{target}, which is currently an unresolved
\textbf{ELEMENT-INAME} or string. If \textbf{*defer-unknown-connections*}
is \textbf{nil}, this instead signals an error. Called internally by
\textbf{connect-wire}; application code normally does not call it
directly.

\textbf{dump-element} (e stream) \emph{{[}FUNCTION{]}}\\
Write a single Lisp form to \textbf{stream} that, when later read by
\textbf{load-kb}, will recreate element \textbf{e} and its wires,
properties, and English names. Used by \textbf{checkpoint-kb} and
\textbf{checkpoint-new}; application code rarely calls it directly.

\paragraph{Manual Checkpointing of the
KB}\label{manual-checkpointing-of-the-kb}

When a KB is growing interactively, we probably want to checkpoint it
from time to time so that knowledge won't be lost if the Scone process
dies or is killed. One way to checkpoint is to save the entire Lisp core
image, but that may not be practical because the core-image may be
multiple gigabytes in size. The functions in this section provide an
alternative.

\textbf{checkpoint-kb} (\&optional filename) \emph{{[}FUNCTION{]}\\
}Dumps all the elements in the current KB into a Lisp (i,e. text) file
with name \textbf{filename}. Any missing components of the filename are
filled in from \textbf{*default-kb-pathname*}. If not supplied, the
filename defaults to \textbf{"checkpoint"}. This does not alter the
running Scone, so it is safe (though perhaps not efficient) to
checkpoint often. The elements created by the \textbf{bootstrap.lisp}
file are not dumped.

To create a clone of the checkpointed Scone process, start a new Scone
process (with \textbf{bootstrap.lisp} already loaded) and then use
\textbf{load-kb} to load the checkpoint file.

\textbf{checkpoint-new} (\&optional filename) \emph{{[}FUNCTION{]}}\\
Like \textbf{checkpoint-kb}, but dumps only the elements created since
the completion of the most recent file load by \textbf{load-kb}. The
default output filename is \textbf{"checkpoint-new"}.

\begin{quote}
To create a clone of the checkpointed Scone process, start a new Scone,
load the same files as in the original Scone, and then load the
checkpoint file. At the start of the checkpoint file, there is a test to
ensure that the expected files are in fact present in the Lisp. (This
checks only for files of the correct filename, not for identical
versions.)
\end{quote}

\textbf{check-loaded-files} (file-list) \emph{{[}FUNCTION{]}}\\
Check whether all the files in \textbf{file-list} are currently loaded
in Scone. Check only the filename itself, not the directory or
extension. Don\textquotesingle t worry about load order or any extra
files that have been loaded. If we fail the test, signal a continuable
error.

\paragraph{Incremental Logging of the
KB}\label{incremental-logging-of-the-kb}

For some applications, occasional checkpointing is not sufficient -- we
want to create a non-volatile log showing \emph{every} change to the KB,
so that we don't lose any knowledge at all if the system goes down
unexpectedly (except perhaps for any KB change that was in progress at
the time of the failure). This logging makes it slower to add new
knowledge to the KB, but in many applications we only care about
inference speed, and modification speed is of less concern.\footnote{A
  continuously streaming log could be made much more efficient if
  implemented at a deep system level, but we will not attempt that for
  now.}

\textbf{start-kb-logging} (\&optional filename) \emph{{[}FUNCTION{]}}\\
KB-logging allows the user to recover from an unexpected termination of
the Scone process. Like \textbf{checkpoint-new}, the
\textbf{start-kb-logging} function first writes out to a text file any
elements created since completion of the most recent \textbf{load-kb}.
It then arranges for any subsequent KB changes to be logged in the same
file. This is implemented by opening a stream bound to the log file and
placing the stream in the global variable \textbf{*kb-logging-stream*},
which controls logging.

The default name for the log file is \textbf{"kb-log"}. Any missing
components of the pathname are filled in from
\textbf{*default-kb-pathname*}.

\begin{quote}
If the Lisp process dies unexpectedly, the user can now recover by
starting a new Scone, loading the same files that were loaded before
logging began, and then calling \textbf{load-kb} on the log file,
ignoring any fragmentary Lisp expression that may be found at the end of
the log. As was the case with \textbf{checkpoint-new}, the log file
begins with a form that verifies that the expected KB files are already
present in the Lisp.
\end{quote}

\textbf{end-kb-logging} nil \emph{{[}FUNCTION{]}\\
}Stops KB logging. Finishes all output and closes the log file.

\textbf{*kb-logging-stream*} \emph{{[}VARIABLE{]}}\\
If non-nil, every KB alteration is reflected by an entry to this stream,
so that the resulting file can be read back in with RELOAD-KB.

\textbf{kb-log} (\&rest args) \emph{{[}FUNCTION{]}}\\
This is the function that actually writes to the KB log. It is called by
any Scone-internal functions that actually modify the KB. Syntax is just
like the Common Lisp \textbf{format} with no stream arg. If
\textbf{*kb-logging-stream*} is non-null, execute the format call,
writing the resulting string out to the KB logging stream; else, do
nothing.

\section{Scone Elements}\label{scone-elements}

\subsection{Structure of Scone
Elements}\label{structure-of-scone-elements}

The Scone knowledge base is best thought of as a collection of
interconnected \emph{elements}. Each element is represented in the Scone
software as a data structure -- a Common Lisp "defstruct".\footnote{Scone
  elements are not currently implemented as objects in the Common Lisp
  Object System (CLOS) because element structures are heavily used in
  inner loops of the Scone engine, and some CLOS implementations are
  inefficient. We may revisit that decision in future releases of Scone.}
An element in Scone can be a \emph{node,} a \emph{link}, or a
\emph{relation}.\footnote{At present, we use a single structure type to
  represent all types of Scone elements. That is a simple solution, but
  is a bit wasteful, since a given element-type will not use all the
  wires and other fields in this structure, though only a relatively
  small fraction of the allocated storage is wasted. We may revisit this
  decision later, and use multiple structure types to more closely fit
  the needs of each element.}

A \emph{node} represents a conceptual entity such as ``George
Washington'' or ``the typical elephant'' or ``the mother of Clyde''.
Note that the node represents a specific \emph{concept} or
\emph{meaning} in each case, not a word or word-definition. The
connection between Scone elements and natural-language words or phrases
is a many-to-many association It is the job of external software (an
English-language front-end to Scone -- under construction) to resolve
any ambiguity.

A \emph{link} represents a relation or statement of some kind, such as
``Clyde is an elephant'' or ``The wife of George is Martha'' or ``Every
elephant hates P. T. Barnum.''

A \emph{split} is a special kind of link with an unusual structure: it
connects to any number of elements rather than just two or three. It is
used to state that these connected elements -- types and individuals --
are mutually distinct and disjoint.

A \emph{relation} is an element that represents the definition of a
relation in Scone. The set of relations is open-ended. Once it has been
created, the relation can be instantiated for specific cases by creating
statement links. A relation is in some respects like a type-node: it can
have other relations below it in the is-a hierarchy -- effectively these
are sub-types of the relation -- and statement-links serve as the
instances.

A Scone knowledge base represents its long-term knowledge as a set of
elements and the pattern of interconnections among these elements. In
general, nodes are connected by links. Each link has a few distinct
\emph{wires} by which it can be connected to other elements in the
knowledge base. These are designated \emph{A-wire}, \emph{B-wire},
\emph{C-wire}, \emph{parent-wire}, \emph{context-wire}, and (for
split-links only) any number of \emph{split-wires}.

The A and B wires are used by links to indicate the two elements that
the link is saying something about. For example, if a link L represents
the statement ``Clyde hates Jumbo'', the A wire of element L would go to
\textbf{\{Clyde\}}, the B wire to \textbf{\{Jumbo\}}, the parent wire to
\textbf{\{hates\}}, and the context wire to the currently-active
context, which is always held in the Lisp variable \textbf{*context*}.

Some relations and statements in Scone also use the \emph{C-wire}. This
allows us to easily represent trinary relations and statements, such as
"the distance from A to B is C".

In general, the \emph{parent wire} of a link is used to indicate one
superior class of this link in the type hierarchy. In other words, the
parent wire tells us what kind of link this is. (A few particularly
important link types, such as \emph{is-a}, \emph{eq}, and \emph{cancel},
are represented as special built-in element types, and are directly
recognizable by the software.) The \emph{context wire} is used for links
to indicate the context (represented by a node) in which a given
statement is considered to be true.

Every element has a \emph{terminal point} to which any number of
incoming wires from other elements may be attached.

We think of a link as connecting two or three nodes (A and B and maybe
C), and representing some statement about the relationship between these
nodes: ``A is a B'' or ``A hates B'' or whatever. However, the link also
is an entity in the knowledge base, representing the statement itself,
and we can make statements (sometimes referred to as
``meta-statements'') about this statement: who told us this piece of
information, how strongly we believe it, whether exceptions are
possible, and so on. So each link element really has two parts: the five
wires, which are used to implement the statement the link is making,
plus the marker bits and terminal point of a node that represents the
statement itself. We sometimes refer to this built-in node as the
\emph{handle node} of the link. (Note that in many cases it is
preferable to put meta-information on a context node containing a
collection of other nodes and links, rather than on the handle node of
an individual link.)

Here is another complication: We think of a node as representing an
entity, about which we can make various statements by attaching
link-wires to it. But every node (except the top-level
\textbf{\{thing\}} node) will be a member of at least one superior
class. So to reduce the total number of elements and to speed up some
common inferences, we equip each node with a \emph{parent wire} that
serves as a virtual is-a link to one superior class. So the
\textbf{\{George Washington\}} individual-node may have a parent-wire
going to the \textbf{\{person\}} type-node, rather than an is-a link to
\textbf{\{person\}}.

While this use of a node's parent-wire significantly reduces the number
of elements required in a Scone knowledge base, there is a cost: because
a parent-wire is not a full-fledged link with an integral handle node,
we cannot attach meta-information to this virtual is-a link or tie it to
a context. If it becomes necessary to do that, we must detach the parent
wire and, replace it with a full-fledged is-a link. This is handled by
the \textbf{convert-parent-wire-to-link} function.

\subsubsection{Functions and Variables}\label{functions-and-variables-1}

\textbf{*n-elements*} \emph{{[}VARIABLE{]}}\\
The number of elements currently in use in the KB.

\textbf{*first-element*} \emph{{[}VARIABLE{]}}\\
First element in the chain of all elements in the KB.

\textbf{*last-element*} \emph{{[}VARIABLE{]}}\\
Last element in the chain of all elements in the KB.

\textbf{do-elements} {[}Complex Macro Body{]} \emph{{[}MACRO{]}}\\
This macro iterates over the set of all elements. A call to this macro
looks like this:\\
\strut \\
\textbf{(do-elements (var) \ldots{} body forms \ldots)\\
} or\\
\textbf{(do-elements (var after) \ldots{} body forms \ldots)\\
\strut \\
}Each element in turn is bound to \textbf{var} and then the body is
executed. Returns \textbf{nil}. If \textbf{after} is supplied, it must
be a form that evaluates to an element. In this case we execute the body
only for elements created after the \textbf{after} element.

\textbf{next-element} (e) \emph{{[}FUNCTION{]}}\\
Given an element \textbf{e}, get the next one in the chain.

\textbf{previous-element} (e) \emph{{[}FUNCTION{]}}\\
Given an element \textbf{e}, get the previous element in the chain.
Note: This is inefficient, since we don\textquotesingle t keep
back-pointers in this chain. It is used in infrequent tasks such as
removing an element from the KB.

\textbf{a-element} (e) \emph{{[}FUNCTION{]}}\\
\textbf{b-element} (e) \emph{{[}FUNCTION{]}}\\
\textbf{c-element} (e) \emph{{[}FUNCTION{]}}\\
\textbf{parent-element} (e) \emph{{[}FUNCTION{]}}

\textbf{context-element} (e) \emph{{[}FUNCTION{]}}\\
Given an element \textbf{e}, return the element connected to the
specified wire of \textbf{e} (the A, B, C, parent, or context wire), or
\textbf{nil} if that wire is not connected.

\textbf{split-elements} (e) \emph{{[}FUNCTION{]}}\\
Given a split element \textbf{e}, return the list of elements connected
to the split wires of \textbf{e}, or \textbf{nil} if there are none.

\textbf{incoming-a-elements} (e) \emph{{[}FUNCTION{]}}\\
\textbf{incoming-b-elements} (e) \emph{{[}FUNCTION{]}}\\
\textbf{incoming-c-elements} (e) \emph{{[}FUNCTION{]}}\\
\textbf{incoming-parent-elements} (e) \emph{{[}FUNCTION{]}\\
}\textbf{incoming-context-elements} (e) \emph{{[}FUNCTION{]}}

\textbf{incoming-split-elements} (e) \emph{{[}FUNCTION{]}}\\
Given an element \textbf{e}, return a list of all elements whose
specified wire (the A, B, C, parent, context, or split wire) is
connected to \textbf{e}'s terminal point. Return \textbf{nil} if there
are none.

\textbf{connect-wire} (wire-name from to \&optional may-defer
already-logged) \emph{{[}FUNCTION{]}}\\
Connect the wire designated by \textbf{wire-name} (one of \textbf{:a},
\textbf{:b}, \textbf{:c}, \textbf{:parent}, \textbf{:context}, or
\textbf{:split}) of element \textbf{from} to element \textbf{to}. If
\textbf{from} already has that wire connected, it is disconnected
first. If \textbf{may-defer} is T and \textbf{to} is not yet a defined
element, the connection is deferred via \textbf{defer-connection} and
resolved later by \textbf{process-deferred-connections}. When the wire
is \textbf{:context} and \textbf{to} is active, \textbf{from} is
activated as well. The change is recorded in the KB log unless
\textbf{already-logged} is T.

\textbf{disconnect-wire} (wire-name from \&optional already-logged)
\emph{{[}FUNCTION{]}}\\
Disconnect the wire designated by \textbf{wire-name} (\textbf{:a},
\textbf{:b}, \textbf{:c}, \textbf{:parent}, or \textbf{:context}) of
element \textbf{from}. If that wire is not currently connected, do
nothing. The change is logged unless \textbf{already-logged} is T. To
disconnect a single \textbf{:split} wire, use
\textbf{disconnect-split-wire}; to disconnect all of them at once, use
\textbf{disconnect-split-wires}.

\textbf{disconnect-split-wires} (from \&optional already-logged)
\emph{{[}FUNCTION{]}}\\
Disconnect every \textbf{:split} wire leaving element \textbf{from}.
Logs the change unless \textbf{already-logged} is T.

\textbf{convert-parent-wire-to-link} (e \&key :context)
\emph{{[}FUNCTION{]}\\
}The parent-wire of element \textbf{e} is disconnected and replaced by
an equivalent is-a link, which can then be modified or tied to a
context. If the user does not supply a \textbf{:context} for the new
is-a link, use the current \textbf{*context*}. Returns the new is-a
link.

\subsection{\texorpdfstring{ Element
Properties}{ Element Properties}}\label{element-properties}

Every Scone element has a \emph{property list}, storing attribute-value
pairs. The set of attributes is open-ended. Each attribute is
represented by a Common Lisp keyword. Some properties are defined and
used by the Scone implementation, while others may be defined by users
or by software external to Scone. These properties provide an
inexpensive way to store meta-information about an element in situations
where the information is static and where high-speed access is not
required; information that may be cancelled or altered, or that may
change from one context to another, is better represented using
full-fledged Scone elements.

Since adding, removing, or changing a property is a modification to the
KB, we normally log this change. (See the \textbf{start-kb-logging}
function.) The \textbf{already-logged} argument, if T, inhibits this
logging.

\subsubsection{Functions}\label{functions}

\textbf{get-element-property} (e property) \emph{{[}MACRO{]}}\\
Find and return the specified property of element \textbf{e}, or
\textbf{nil} if this property is not present. The \textbf{property}
argument is typically a Lisp keyword.

\textbf{set-element-property} (e property \&optional value
already-logged) \emph{{[}FUNCTION{]}}\\
Set the specified property of element \textbf{e} to \textbf{value}. If
\textbf{value} is not supplied, the default value is \textbf{t}. Returns
\textbf{value}.

\textbf{clear-element-property} (e property \&optional already-logged)
\emph{{[}FUNCTION{]}}\\
Remove the specified \textbf{property} of element \textbf{e}. If
\textbf{e} doesn\textquotesingle t have this property, do nothing.

\textbf{push-element-property} (e property value \&optional
already-logged) \emph{{[}FUNCTION{]}}\\
The \textbf{property} of \textbf{e} should be a list. Push
\textbf{value} onto this list. Create the property if it
doesn\textquotesingle t already exist.

\section{\texorpdfstring{ Referring to Scone
Elements}{ Referring to Scone Elements}}\label{referring-to-scone-elements}

Many of the Scone functions described in this document take one or more
\emph{elements} as arguments. Within the running Scone system, each
element is represented as a Common Lisp structure (sometimes called a
``defstruct''). If you have a pointer to one of these element-objects,
perhaps returned by \textbf{lookup-element} or some other Scone
function, you can supply this directly wherever an element is
required.\footnote{In Lisp, almost all data objects are allocated in
  dynamic storage, and we pass around pointers to these objects. So it
  is common in Lisp culture to say that we are ``passing Scone element X
  to function Y'', when in fact we are really passing a pointer to the
  heap-allocated structure representing Scone element X. If all pointers
  to the element-X structure are dropped, the Lisp garbage collector
  will eventually reclaim the storage. But note that a chain of pointers
  runs through all the Scone elements, so no element structure will be
  reclaimed unless we explicitly snip it out of this chain, perhaps with
  \textbf{remove-element}.}

Alternatively, you can refer to an element by its name. There are two
kinds of element names in Scone: \emph{internal names} (sometimes called
\emph{inames}) and \emph{external names} (sometimes called ``English
names'', though other natural languages will eventually be supported).

An internal name such as \textbf{\{George Washington\}} refers uniquely
and unambiguously to a specific Scone element. External names are not so
constrained: an element may have many external names, and an external
name such as \textbf{"mouse"} may refer to several distinct Scone
elements -- in this case, one representing the animal mouse and one
representing a computer pointing device. A knowledge-base file can
contain both internal and external names, but if a reference is made to
an external name with more than one possible meaning, the system may
interactively request a clarification while the file is being loaded. So
it is generally best to use internal names within KB files.

\subsection{Element Internal Names}\label{element-internal-names}

Internal names are surrounded by ``curly braces''. So, for example, the
internal name of the node at the top of the type hierarchy can be
referred to as \textbf{\{thing\}}. This curly-brace notation can be used
wherever a Scone element is required.

An internal name is an arbitrary character string that serves as the
label for a specific element. We tend to use names like
\textbf{\{elephant\}} that will be meaningful to a human reader, but
that is just for convenience in developing and understanding the
system's knowledge base. As far as Scone is concerned, the internal name
for the \textbf{\{elephant\}} node could just as well be
\textbf{\{elephant (the animal)\}}, \textbf{\{P123456\}},
\textbf{\{xyzzy\}}, or \textbf{\{aardvark\}} -- as long as the chosen
name is unique.

These iname strings are matched in a case-insensitive way; that is,
\textbf{\{Clyde\}}, \textbf{\{clyde\}}, and \textbf{\{CLYDE\}} all refer
to the same element. However, Scone remembers the exact string,
including case, that was used when an element was first created, and
that string is used whenever the system must print out a reference to
the element.

Internal name strings may contain spaces and other punctuation
characters, but not the colon character, which is reserved for another
use (see below). So, for example,\\
\textbf{\{mother of Clyde\}} is a legal internal name, and names with
spaces are indeed fairly common.

Internal names are organized into \emph{namespaces}. Within its
namespace, an internal name must be unique. That is, it refers to only
one element. Any attempt to create a second node with the same iname in
a given namespace is an error.

Namespaces are designated by Common Lisp strings such as
\textbf{"common"} or \textbf{"astronomy"}. These \emph{namespace names}
must be globally unique so that accidental collisions do not occur if we
try to combine separately developed knowledge bases. An informal
registry of Scone namespace names currently in use is maintained by the
author. Contact the author if you want to register a new global
namespace. (Scone may eventually have to move to a more complex
universal namespace scheme, of the kind found in XML-based knowledge
representations, but we want to delay that move as long as possible --
the resulting names are long and confusing for human readers.)

To designate an internal name without danger of ambiguity we use its
\emph{full name}. This is a Common Lisp string divided into two parts by
a colon. The part of the string before the colon designates the
namespace, while the part after the colon is the unique internal name
within that namespace. So \textbf{\{biology:mouse\}} and
\textbf{\{computer:mouse\}} may co-exist without creating a conflict.
There may be only one colon in an internal name string.

There is a function called \textbf{in-namespace} that takes a string as
an argument. That string is converted to a namespace object, which
becomes the value of the \textbf{*namespace*} variable. Whenever the
system is given an iname without a colon in it (referred to as a
\emph{short name}), the value of the \textbf{*namespace*} variable is
used as the default namespace. Typically a KB file will have an
\textbf{in-namespace} form at the start and will use short-names
internally, except when it is necessary to refer to an iname in another
namespace.

A namespace may optionally \emph{include} one other namespace, which may
include another, and so on, forming a chain. If an attempt is made to
look up an internal name in a given namespace and the iname is not found
there, the system will look in the included namespace (if any) and then
in its included namespace, and so on until the end of the chain is
reached. If the iname is not found anywhere in this chain of namespaces,
the name is undefined.

If we create a new iname, it will be placed directly into the current
\textbf{*namespace*}, but an error will be signaled if the name
conflicts with the same namestring in any included namespace. This
namespace-inclusion mechanism is often used to include the
\textbf{"common"} namespace in some more specialized namespace that the
user is creating, so that elements like \textbf{\{common:thing\}} and
\textbf{\{common:person\}} can be referred to simply as
\textbf{\{thing\}} and \textbf{\{person\}}.

Links normally do not have meaningful internal names. However, if you
believe that you will want to refer to a specific link, you can give it
an internal name via the \textbf{:iname} keyword argument of the
``new-whatever'' function that creates the link. The
\textbf{*generate-long-element-names*} control variable, if \textbf{t},
tells the system to create longer, more meaningful names for links. That
requires more space, but can be useful for development and debugging.

For elements that represent primitive Lisp data types, such as numbers,
functions, or strings, the Lisp object itself is used as the iname and
there is never a namespace designator. (These objects are kept in a
special universal namespace.) So \textbf{\{123\}} designates an element
representing an integer, and \textbf{\{12.34\}} represents a real
number. The Scone element \textbf{\{\#\textquotesingle query-user\}}
represents a Lisp function named \textbf{query-user}.

A set of brackets containing a quoted string, such as
\textbf{\{"http://www.cmu.edu"\}}, is a string element, representing the
Lisp string itself. We use these string elements to represent words,
names, URLs, command-line commands, and so on. Note that any colon
within a string element is part of the string, not a namespace
designator.

\subsubsection{Functions and Variables}\label{functions-and-variables-2}

\textbf{*namespace*} \emph{{[}VARIABLE{]}}\\
The structure representing the current default namespace.

\textbf{*namespaces*} \emph{{[}VARIABLE{]}}\\
A hash table containing an entry for each namespace currently loaded in
the KB. Associates a string (the name of the namespace) with a namespace
structure.

\textbf{list-namespaces} nil \emph{{[}FUNCTION{]}}\\
Return a list of all the Scone namespaces that are currently loaded.

\textbf{lookup-element} (name \&key :syntax-tags) \emph{{[}FUNCTION{]}\\
}This takes an element-designator \textbf{name} (an element object, an
internal name in curly braces, or an external name string) and returns
the corresponding element object. If the argument does not correspond to
an existing element in the KB, return \textbf{nil}.

The optional \textbf{:syntax-tag} argument, if present, is used in the
resolution of external (English) name strings to element objects. (See
the discussion of the \textbf{disambiguate} function in the English
Names section of this manual). If \textbf{lookup-element} is given a
string argument and successfully finds an element, it returns the
element\textquotesingle s syntax-tag as the second return value.

\textbf{lookup-element-test} (name \&key :syntax-tags)
\emph{{[}FUNCTION{]}}\\
Like \textbf{lookup-element}, but if the named element does not (yet)
exist, signal an error.

\textbf{lookup-element-or-defer} (name \&key :syntax-tags)
\emph{{[}FUNCTION{]}}

Like \textbf{lookup-element}, but if the named element does not (yet)
exist, return the argument unchanged.

\textbf{in-namespace} (namespace-name \&key :include)
\emph{{[}FUNCTION{]}}\\
The namespace designated by the \textbf{namespace-name} string becomes
the new default namespace (the value of the \textbf{*namespace*}
variable). If the specified namespace does not already exist, it is
created. If the \textbf{:include} argument is supplied, this indicates
an existing namespace that is included in the new one, if a new one is
actually created. The \textbf{:include} argument may be either a
namespace object or a string designating a namespace.

\textbf{make-namespace} (namespace-name \&key :include)
\emph{{[}FUNCTION{]}}\\
Create a new, empty namespace with the given \textbf{namespace-name}
string and register it. The optional \textbf{:include} argument, if
supplied, specifies a namespace (by object or string) whose entries
will be visible from the new namespace. This is the lower-level
constructor used by \textbf{in-namespace}; most user code should call
\textbf{in-namespace} instead.

\textbf{get-namespace} (namespace-name \&key :include)
\emph{{[}FUNCTION{]}}\\
Return the namespace object registered under \textbf{namespace-name},
creating it (with the given \textbf{:include}) if it does not already
exist. The lookup is case-insensitive.

\textbf{extract-namespace} (string) \emph{{[}FUNCTION{]}}\\
Given a string, return the portion before the first colon, or
\textbf{nil} if there is no colon. Used to separate the namespace part
of a qualified internal name from the element-name part.

\textbf{read-element} (stream char) \emph{{[}FUNCTION{]}}\\
The Common Lisp reader macro for the \textbf{\{} character. When the
reader sees an open curly brace, it calls this function to read the
contents of the braces and build an \textbf{ELEMENT-INAME} structure
of type \textbf{:ELEMENT}, \textbf{:STRING}, \textbf{:NUMBER}, or
\textbf{:FUNCTION}, which is later resolved to a real element by
\textbf{lookup-element}. Application code rarely calls this directly;
it is documented here because it is what gives Scone its characteristic
\textbf{\{curly-brace\}} syntax.

\textbf{lookup-element-iname} (name) \emph{{[}FUNCTION{]}}\\
Specialized form of \textbf{lookup-element} for the case where the
argument is already known to be an \textbf{ELEMENT-INAME} structure
(for instance, one just produced by \textbf{read-element}). Handles the
four iname types -- \textbf{:ELEMENT}, \textbf{:STRING},
\textbf{:NUMBER}, \textbf{:FUNCTION} -- and resolves explicit
\textbf{namespace:name} prefixes.

\textbf{element-name} (e) \emph{{[}FUNCTION{]}}\\
Return a printable string for element \textbf{e} in the usual
curly-brace form, optionally including the namespace and colon. The
variable \textbf{*print-namespace-in-elements*} controls whether the
namespace is included: \textbf{t} always includes it, \textbf{nil}
never does, and \textbf{:maybe} includes it only when the element is
from a namespace other than the current one.

\textbf{iname} (x) \emph{{[}FUNCTION{]}}\\
Return the bare internal name of \textbf{x}, which may be an element
or an \textbf{ELEMENT-INAME} structure. The result is normally a
string; for primitive nodes (numbers, functions, etc.) it may be the
underlying Lisp object. Unlike \textbf{element-name}, this does not
include curly braces or namespace information.

\textbf{print-element-iname} (e stream depth)
\emph{{[}FUNCTION{]}}\\
The print-function for \textbf{ELEMENT-INAME} structures. Writes
\textbf{e} to \textbf{stream} in the usual \textbf{\{name\}} form.
Installed automatically by the \textbf{ELEMENT-INAME} \textbf{defstruct};
normally not called directly.

\subsection{Element External (English)
Names}\label{element-external-english-names}

A long-term goal for Scone is to interface it to a natural-language
input-output system so that users can augment or query the knowledge
base in English (or in the human language of their choice). At present,
the system has only a simple facility for associating words or phrases
(represented as Common Lisp strings) with various elements in the
knowledge base. We refer to these words and phrases collectively as
\emph{external names} to distinguish them from the internal
element-names described in the previous section.

While most of the machinery for natural-language input and output is
external to the Scone inference engine, it makes sense (for now, at
least) to maintain the dictionary mapping names to elements and elements
to names as part of the core Scone system. When a KB file containing new
concepts is loaded into the knowledge base, the external names for these
concepts are loaded at the same time, as part of the information in this
file.

The mapping between external names and Scone elements is many-to-many.
Scone provides a simple interface by which an external natural-language
system can ask for all the Scone elements (if any) that correspond to a
given name. Similarly, the user can ask for a list of all the names (if
any) that correspond to a given element.

Note that the name strings are matched exactly, except that the match is
case-insensitive. At present there is no spelling correction or ``find
an approximate match'' capability. In general, external name strings
should be entered in lower-case, except in the case of proper nouns like
``Obama''. There should be no whitespace surrounding the word or phrase.
In multiple-word strings (representing idiomatic phrases such as
\textbf{``kick the bucket''}) the words should be separated by a single
space.

Each mapping of a name-string to an element has an associated
\emph{syntax-tag}. This is a Lisp keyword that corresponds roughly to
the part of speech associated with the name-string. This allows the NL
machinery to return only those Scone meanings that correspond to the
part of speech it has found -- for example, the action ``fall'' rather
than a noun representing a season.

\emph{Note: This syntax-tag mechanism is something of a place-holder for
now. Once the natural-language interface to Scone is further developed,
we will have a much better idea of what syntactic machinery really is
needed. We will probably end up with a complex hierarchy of word-types
and grammatical roles, represented in Scone itself, but more or less
distinct from the conceptual knowledge that the linguistic forms refer
to.}

For now, the following syntax-tags are defined:

\textbf{:noun\\
}Either a common name or a proper name for the entities represented by
the element. "Clyde" or "elephant".

\textbf{:adj\\
}An adjective like "blue", which represents a unary predicate of some
kind, and is generally associated with a Scone type-node. To describe a
type, you normally combine this adjective with some noun labeling a
supertype of the type-node in question. So we would normally say "a blue
object" or "a blue automobile" rather than just "a blue".

\textbf{:adj-noun\\
}A word like "solid" that can be used either as a noun or as an
adjective. These are very common in English.

\textbf{:role\\
}A word like "population" that designates a role under some type node.
Normally this used with "of" or a possessive form: "the population of
India" or "India\textquotesingle s population".

\textbf{:inverse-role\\
}Name given to the inverse of a role relation. For example, if the role
name is ``parent'', the inverse-role may be ``child''. So if we are told
that "John is the child of Mary", the natural-language machinery can
convert this to "Mary is the parent of John".

\textbf{:relation\\
}A word like "employs" or "hates" or a phrase like "taller than"
representing some relation that is (grammatically) not a role.

\textbf{:inverse-relation\\
}Name given to the inverse of a relation. For example, if the relation
name is ``taller than'', the inverse-relation may be ``shorter than''.
So if we are told that "John shorter than Mary", the natural-language
machinery can convert this to "Mary is taller than John".

\textbf{:action\\
}An action verb such as "hit".

\subsubsection{Functions and Variables}\label{functions-and-variables-3}

\textbf{*legal-syntax-tags*} \emph{{[}VARIABLE{]}}\\
Each external name has an associated syntax-tag. The tags currently
defined are stored in this variable.

\textbf{english} (element \&rest r) \emph{{[}FUNCTION{]}}\\
Define one or more English names for the specified \textbf{element}. The
remaining arguments after element are strings or syntax tags (keywords).
Initially, we assume that each string becomes an English name of element
with syntax-tag \textbf{:noun}, but the tag changes in a sticky way
whenever a different syntax-tag is encountered. If the special
syntax-tag \textbf{:iname} is encountered, that indicates that the
element\textquotesingle s internal name is used as one of the English
names. However, if the argument list contains \textbf{:no-iname}, that
takes precedence and any \textbf{:iname} in the \textbf{r} list is
ignored.\\
\strut \\
For example, the form\\
\textbf{(english \{swamp\} :iname "wetland" "marsh" :adj "soggy"
"marshy")}\\
says that the element \textbf{\{swamp\}} has English names
\textbf{"swamp"}, \textbf{"wetland"}, and \textbf{"marsh"} (all nouns),
plus the adjectives \textbf{"soggy"} and \textbf{"marshy"} (as in "soggy
area" or "marshy area").

\textbf{lookup-definitions} (string \&optional syntax-tags)
\emph{{[}FUNCTION{]}}\\
Returns a list containing all the meanings of the \textbf{string}
argument (if any). Each meaning in the list will be of the form
\textbf{(element . syntax-tag)}. If the \textbf{syntax-tag} argument is
supplied, construct and return a list containing only meanings with the
specified syntax tag. This function works only for external name strings
and is called by the other functions of the lookup-element family.

\textbf{mark-named-elements} (string m \&optional syntax-tag)
\emph{{[}FUNCTION{]}}\\
All the elements associated with the name \textbf{string} are marked
with marker \textbf{m}. If \textbf{syntax-tag} is supplied, only
meanings of the specified type are marked.

\textbf{get-english-names} (element \&optional syntax-tag)
\emph{{[}FUNCTION{]}}\\
Returns a list of all the names associated with \textbf{element} (if
any). Each name will be of the form \textbf{(string . syntax-tag)}. If a
\textbf{syntax-tag} argument is supplied, construct and return a list
containing only names of the specified type.

\textbf{disambiguate} (name definition-list \&optional syntax-tags)
\emph{{[}FUNCTION{]}}\\
Given a list of possible definitions (element.tag pairs) for an English
word, try to disambiguate it based on the value of
*disambiguate-policy*. If :ERROR, just signal an error. If :ASK, ask the
user. If :HEURISTIC, use SYNTAX-TAGS to choose the best definition, or
if that fails just pick one.

\textbf{list-synonyms} (string \&optional syntax-tag)
\emph{{[}FUNCTION{]}}

\textbf{show-synonyms} (string \&optional syntax-tag)
\emph{{[}FUNCTION{]}}\\
The \textbf{list-synonyms} function finds all the elements that are
registered as a meaning of \textbf{string} in the KB. For each of these
it returns a list, each element of which has a Scone element first,
followed by any number of alternative names for that element. If
\textbf{syntax-tag} is non-null, we look only at elements whose
connection to \textbf{string} is of this type. The
\textbf{show-synonyms} function is similar, but it prints out the
synonyms it finds instead of returning a list.

\textbf{list-dictionary-entries} nil \emph{{[}FUNCTION{]}}\\
List all entries in the dictionary. This is primarily for debugging.

See also the functions \textbf{lookup-element},
\textbf{lookup-element-predicate}, and \textbf{lookup-element-test} in
the previous section. These can be used to convert a string (and
optional syntax-tag) to a single Scone element. If the string is
ambiguous -- that is, if it is associated with more than one Scone
element via the given syntax-tag -- these functions call
\textbf{disambiguate}, which asks the user to pick one of the possible
meanings.

\section{\texorpdfstring{Markers }{Markers }}\label{markers}

Every Scone element has a set of bits representing permanent state
information about that element. These bits are known as the element's
\emph{flag bits}. Some of the flag bits are used to encode the type of
the node or link: individual vs. type node, ``is a'' versus ``eq'' link,
and so on. Other flag bits encode properties of the element, such as
whether a statement link is an instance of a transitive relation. As a
rule, flag bits are intended for internal use by the Scone engine, and
are not directly manipulated by users of the Scone system, so we will
not describe them in detail in this manual.

Every Scone element also has storage for a number of \emph{marker bits},
which are set and cleared by the Scone engine as part of various search
and inference operations. Operations on marker bits generally refer to a
specific marker using a numerical index, starting at 0. The constant
\textbf{n-markers} indicates the maximum number of marker bits a Scone
element can hold: in the current implementation for 32-bit SBCL,
\textbf{n-markers} is 28, so legal markers are designated by integers 0
through 27. On a 64-bit Common Lisp, there are more markers available --
check \textbf{n-markers} to see how many you have.

Markers are normally allocated and used in pairs: a
\emph{positive-marker}, used to mark some set of elements, and an
associated \emph{cancel-marker}. The cancel-marker is placed on elements
\emph{definitely not} in the set marked with the positive marker and
also on links that are cancelled in the course of marking this set of
elements. We speak of a marker and its cancel marker together as a
\emph{marker-pair}. Effectively, these pairs implement a 3-valued logic:
yes, definitely no, and maybe or ``don't know''.

The number of possible marker pairs, as well as the number of positive
markers, is indicated by the constant \textbf{n-marker-pairs}, which
will be half of \textbf{n-markers} (rounded down). In the 32-bit SBCL
implementation \textbf{n-marker-pairs} is 14. So markers numbered 0-13
are positive markers, and markers 14-27 are the corresponding cancel
markers. When given a positive marker \textbf{m}, the function
\textbf{(cancel-marker m)} will return the corresponding cancel marker.

Two of these marker-pairs are currently reserved for use by the Scone
engine. The \emph{context marker} (normally marker-pair 0) is used to
mark all the nodes representing the currently active context and its
superiors. The global variable \textbf{*context-marker*} always holds
the index of the context marker, so users should refer to the marker via
that variable. Similarly, the \textbf{*activation-marker*} (normally
marker-pair 1) is used to mark all the elements that are active in the
current context. The remaining marker-pairs are available for search and
inference operations initiated by the user.

Scone provides functions for allocating and freeing markers (or
marker-pairs) for use by various software operations. The variable
\textbf{*available-markers*} indicates the number of marker-pairs
available for allocation at any given time. The \textbf{get-marker}
function selects and reserves one of the available markers and returns
its index; the corresponding cancel-marker is also reserved by this
operation. The \textbf{free-marker} function removes the specified
marker and its cancel-marker from all Scone elements and returns this
marker-pair to the available-markers pool. The
\textbf{with-temp-markers} macro is a convenient way to allocate a
temporary marker-pair, use it while executing any number of operations,
and then to reliably free it.

For a detailed description of how Scone uses these markers for search
and inference, please see the paper "Marker-Passing Inference in the
Scone knowledge-Base System", available from the Scone website:
\url{http://www.cs.cmu.edu/~sef/scone/publications/}.

\subsubsection{Functions and Variables}\label{functions-and-variables-4}

\paragraph{Define the Range of Legal Markers and
Marker-Pairs}\label{define-the-range-of-legal-markers-and-marker-pairs}

\textbf{n-markers} \emph{{[}CONSTANT{]}}\\
The maximum number of marker bits allowed in this Scone engine. This
number is set when the Scone engine is compiled for a particular
hardware architecture, so it cannot be changed dynamically.

\textbf{n-marker-pairs} \emph{{[}CONSTANT{]}}\\
The maximum number of marker pairs.

\textbf{legal-marker?} (m) \emph{{[}MACRO{]}}

\textbf{legal-marker-pair?} (m) \emph{{[}MACRO{]}}\\
Predicates to determine whether a number \textbf{m} represents a legal
marker in the current Scone. The \textbf{legal-marker?} predicate
accepts any integer greater than 0 and less than \textbf{n-markers}; the
\textbf{legal-marker-pair?} predicate accepts only those integers less
than \textbf{n-marker-pairs?}.

\textbf{check-legal-marker} (m) \emph{{[}MACRO{]}}

\textbf{check-legal-marker-pair} (m) \emph{{[}MACRO{]}}\\
Check whether \textbf{m} is a legal user marker. If not, signal an
error.

\textbf{marker-bit} (m) \emph{{[}MACRO{]}}\\
Get the marker bit mask associated with marker \textbf{m}.

\textbf{get-cancel-marker} (m) \emph{{[}MACRO{]}}\\
Get the cancel-marker corresponding to marker \textbf{m}.

\paragraph{Basic Marker Functions}\label{basic-marker-functions}

\textbf{mark} (e m) \emph{{[}FUNCTION{]}}\\
Marks element \textbf{e} with marker \textbf{m}. If \textbf{e} is
already marked with \textbf{m}, do nothing. Return the number of
elements marked with \textbf{m} after this operation.

\textbf{unmark} (e m) \emph{{[}FUNCTION{]}}\\
Removes marker \textbf{m} from element \textbf{e}. If \textbf{e} is not
marked with \textbf{m}, do nothing. Return the number of elements marked
with \textbf{m} after this operation.

\textbf{marker-count} (m) \emph{{[}FUNCTION{]}\\
}Returns the number of elements currently marked with marker \textbf{m}.

\textbf{marker-on?} (e m) \emph{{[}FUNCTION{]}}\\
Predicate to determine whether element \textbf{e} has marker \textbf{m}
turned on.

\textbf{marker-off?} (e m) \emph{{[}FUNCTION{]}}\\
Predicate to determine whether element \textbf{e} has marker \textbf{m}
turned off.

\textbf{clear-marker} (m) \emph{{[}FUNCTION{]}}\\
Clear marker M from all elements, but do not free it.

\textbf{clear-marker-pair} (m) \emph{{[}FUNCTION{]}}\\
Clear marker M and the associated cancel marker from all elements, but
do not free it.

\textbf{clear-all-markers} nil \emph{{[}FUNCTION{]}}\\
Clear and free all markers from all elements in the KB. Note that this
clears the context and activation markers, so you should follow this
with an \textbf{in-context} or \textbf{refresh-context} call.

\textbf{do-marked} (var m \&optional after) \ldots{} body forms \ldots{}
\emph{{[}MACRO{]}\\
}This macro iterates over the set of all elements marked with marker
\textbf{m}. A call to this macro looks like this:\\
\strut \\
\textbf{(do-marked (var m) \ldots{} any number of body forms \ldots)\\
} or\\
\textbf{(do-marked (var m after) \ldots{} any number of body
forms\ldots)\\
\strut \\
}Each \textbf{m}-marked element in turn is bound to \textbf{var} and
then the body is executed. Returns \textbf{nil}. If \textbf{after} is
supplied, it must be a form that evaluates to an element marked with
\textbf{m}. Only iterate over elements after this one in the chain of
\textbf{m}-marked elements. \emph{NOTE: It is OK to mark new elements
with \textbf{m} in the body of this loop, but don\textquotesingle t
unmark any elements.}

\paragraph{Boolean Marker Functions}\label{boolean-marker-functions}

\textbf{mark-boolean} (m must-be-set must-be-clear \&optional flags-set
flags-clear) \emph{{[}FUNCTION{]}}\\
The \textbf{m} argument is a marker. The \textbf{must-be-set} and
\textbf{must-be-clear} arguments are lists of markers. Conceptually, we
scan every element in the KB. If all of the \textbf{must-be-set} markers
are on and none of the \textbf{must-be-clear} markers are on, turn on
marker \textbf{m} for this element. Returns the number of elements
marked with m after this operation.

\begin{quote}
So, for example, \textbf{(mark-boolean 3 \textquotesingle(0 1)
\textquotesingle(2))} places marker 3 on every element that is currently
marked with 0 and 1 and that is not marked with 2.

The \textbf{flags-set} and \textbf{flags-clear} arguments, if supplied,
can control which kinds of elements will potentially be marked. This
option is for internal use, and is not documented here.
\end{quote}

\textbf{unmark-boolean} (m must-be-set must-be-clear \&optional
flags-set flags-clear) \emph{{[}FUNCTION{]}}\\
Like \textbf{mark-boolean}, but clears marker \textbf{m} from the
selected elements.

\textbf{convert-marker} (m1 m2) \emph{{[}FUNCTION{]}}\\
For every element marked with \textbf{m1}, mark it with \textbf{m2} (if
it was not marked with \textbf{m2} already) and clear \textbf{m1} from
that element. This could be done with \textbf{mark-boolean}, but this
specialized version is faster.

\paragraph{Marker Allocation}\label{marker-allocation}

Note: It is a Very Bad Idea to just grab a marker and start using it
without informing the system that this marker is now in use. This can
lead to bugs that are very hard to diagnose. So always use
\textbf{with-markers} or \textbf{get-marker} and \textbf{free-marker}
when you want a marker to play with.

\textbf{*n-available-markers*} \emph{{[}VARIABLE{]}}\\
The number of markers currently available for allocation.

\textbf{get-marker} nil \emph{{[}FUNCTION{]}}\\
Allocate an available marker, returning the integer index. If no
available markers remain, return \textbf{nil}.

\textbf{free-marker} (m) \emph{{[}FUNCTION{]}}\\
Return Marker \textbf{m} to the pool of available markers and clear the
marker. If the marker is already available or is reserved by the system,
do nothing.

\textbf{free-markers} nil \emph{{[}FUNCTION{]}}\\
Return all markers to the free pool and clear them, except for those
such as the context marker that are permanently allocated to the Scone
engine.

\textbf{with-markers} (marker-vars body-form \&optional fail-form)
\emph{{[}MACRO{]}\\
}The \textbf{marker-vars} argument is a list of variables. Each is
locally bound to the index of a freshly allocated marker. We then
execute the body form, returning whatever value or values that form
returns. But before leaving this form, either in the normal way or via
an error or non-local exit, we de-allocate and clear these markers.

\begin{quote}
If there are not enough markers available to fill all the
\textbf{marker-vars}, we don\textquotesingle t allocate any new markers.
Instead we execute the \textbf{fail-form} and return whatever it
returns. The default action, if there is no \textbf{fail-form}, is to
signal an error.

A call to this macro might look like this:\\
\textbf{(with-markers (v1 v2 v3)\\
(progn\\
(some-action-using-marker v1)\\
(another-action-using-markers v1 v2 v3))\\
(deal-with-an-allocation-failure))}
\end{quote}

\textbf{available-marker?} (m) \emph{{[}MACRO{]}}\\
Predicate that returns \textbf{t} if marker \textbf{m} is currently
in the free pool (that is, not allocated to anyone).

\textbf{lock-marker} (m) \emph{{[}FUNCTION{]}}\\
Mark marker \textbf{m} as locked. A locked marker is not returned to
the free pool by \textbf{free-marker} or \textbf{free-markers}. Scone
uses this internally to protect the context and activation markers
(allocated by \textbf{context-marker-setup}), but application code can
also lock a marker it wants to keep for the lifetime of a process.

\textbf{unlock-marker} (m) \emph{{[}FUNCTION{]}}\\
Remove the lock on marker \textbf{m}, allowing it to be freed
normally.

\textbf{locked-marker?} (m) \emph{{[}FUNCTION{]}}\\
Predicate that returns \textbf{t} if marker \textbf{m} is locked.

\textbf{context-marker-setup} \emph{{[}FUNCTION{]}}\\
Allocate and lock the two markers that Scone uses internally for
context and activation tracking. The bootstrap KB calls this once at
startup; it is normally not called by application code. A second call
on an already-initialized process is silently ignored.

\paragraph{Low-level Bit Predicates}\label{low-level-bit-predicates}

These macros operate on raw bit-vector representations of marker sets
-- the \textbf{bits} field of an element, or a mask produced by
\textbf{marker-bit}. They are used internally by the marker scans but
are also available for applications that need to test several markers
at once without looping.

\textbf{all-bits-on?} (bits mask) \emph{{[}MACRO{]}}\\
Returns \textbf{t} if every one-bit in \textbf{mask} is also on in
\textbf{bits}.

\textbf{any-bits-on?} (bits mask) \emph{{[}MACRO{]}}\\
Returns \textbf{t} if at least one bit set in \textbf{mask} is also on
in \textbf{bits}.

\textbf{all-bits-off?} (bits mask) \emph{{[}MACRO{]}}\\
Returns \textbf{t} if no bit set in \textbf{mask} is on in
\textbf{bits}.

\section{Contexts}\label{contexts}

Every Scone element has a context-wire. This is used to tie a node to
the context in which it exists, and it is used to tie a link to the
context is which it is true. A context is simply a Scone node. When a
new context-node is first created, it inherits the contents (nodes and
links) of its superior nodes in the is-a hierarchy. You can then add
additional KB structure in the new context. You can also cancel or
negate nodes and links that would otherwise be inherited.

The care and feeding of contexts, and all the uses to which they can be
put, is a big topic, beyond the scope of this manual. Here we will just
document that functions and variables that implement the context
machinery.

\subsubsection{Functions and Variables}\label{functions-and-variables-5}

\textbf{*context*} \emph{{[}VARIABLE{]}}\\
The element representing the current context..

\textbf{in-context} (e) \emph{{[}FUNCTION{]}\\
}Argument \textbf{e} is an existing element representing a context.
\textbf{e} becomes the value of \textbf{*context*}. Then we propagate
\textbf{*context-marker*} to all the contexts whose contents are
inherited by \textbf{e}, and we propagate \textbf{*activation-marker*}
to all the elements active in this set of contexts. If the
\textbf{*context*} is already \textbf{e}, do nothing. This can be a
relatively expensive operation -- we have chosen to do more work when
changing contexts in order to make within-context operations as fast as
possible.

\textbf{new-context} (iname \&optional parents) \emph{{[}FUNCTION{]}}\\
This is a convenience function for creating a new context node as a
descendant of one or more pre-existing context-nodes. It is actually a
specialized form of \textbf{new-indv} (see below). The \textbf{parents}
argument may either be a single parent node or a list of parents. If
\textbf{parents} is unsupplied or \textbf{nil}, assume that there is a
single parent, \textbf{\{general\}}.

\textbf{refresh-context} nil \emph{{[}FUNCTION{]}}\\
Keep the current value of \textbf{*context*}, but update the context and
activation markers. This is used when we think that the context and
activation markers may somehow have been messed up.

\textbf{*context-marker*} \emph{{[}VARIABLE{]}}\\
Marker permanently assigned to mark the current \textbf{*context*} node
and its superiors.

\textbf{*context-cancel-marker*} \emph{{[}VARIABLE{]}}\\
Cancellation marker corresponding to \textbf{*context-marker*}.

\textbf{*activation-marker*} \emph{{[}VARIABLE{]}}\\
Marker permanently assigned to mark every element active or present in
the current context.

\textbf{*activation-cancel-marker*} \emph{{[}VARIABLE{]}}\\
Cancellation marker corresponding to \textbf{*activation-marker*}.

\section{Adding Elements to the Scone
KB}\label{adding-elements-to-the-scone-kb}

\subsection{Creating New Elements}\label{creating-new-elements}

Each type of Scone element has a ``new'' function that creates it:
\textbf{new-indv}, \textbf{new-type}, \textbf{new-is-a}, and so on.
These new-functions also check whether the element about to be added
conflicts with other information currently in the KB.

All of these new-functions have certain features in common. All will
create at least one new element in the KB, and the new element object
created (or the \emph{primary} new element created, if there are more
than one) is returned as the value of executing the new-function. A
knowledge-base (or KB) file is just a collection of such new-functions,
with perhaps some comments and assorted other functions mixed in.

New-functions that create a node, such as \textbf{new-indv}, generally
take two required arguments, the internal name of the new node and the
parent node. So a typical call might be something like

\textbf{(new-indv \{Clyde\} \{elephant\})}.

Most new-functions that create a link take two required arguments, the
two elements that the link connects, called the A and B arguments. So a
typical call might be something like

\textbf{(new-is-a \{Clyde\} \{elephant\})}.

The new-function for general statements takes three required arguments,
as in

\textbf{(new-statement \{Clyde\} \{hates\} \{Ernie\})}.

In this case, \textbf{\{Clyde\}} is attached to the A-wire of the new
statement link, \textbf{\{Ernie\}} is attached to the B-wire, and
\textbf{\{hates\}} is attached to the parent-wire of the link.

\subsubsection{The :context Keyword
Argument}\label{the-context-keyword-argument}

All new-functions take optional keyword arguments. Normally, all new
elements are created in the current context (as indicated by the current
value of the \textbf{*context*} variable). The new elements'
context-wires are connected to that context node. However, you can
over-ride this for a particular new element by providing a
\textbf{:context} argument pointing to some other context-node.

\subsubsection{\texorpdfstring{\textbf{Internal Names and the :iname
Argument}}{Internal Names and the :iname Argument}}\label{internal-names-and-the-iname-argument}

\textbf{iname} is not a required argument for links, splits, and certain
other elements. The system normally just makes up a suitable new name
for these elements. However, you can provide a specific iname for these
elements if you like via the \textbf{:iname} keyword argument.

Where an \textbf{iname} argument is required, a value of \textbf{nil}
tells Scone to make up a unique name for the new element. For example,
if we say

\textbf{(new-indv nil \{elephant\})}

the new individual node will be given a name like \textbf{\{elephant
(0-27)\}}. Numbers are appended to these made-up-names to ensure that
they are unique.

For convenience, the \textbf{iname} argument can be a quoted string
instead of an internal name in curly braces, but the effect is the same:
that string is converted into an iname, and it must be unique within its
namespace. So \textbf{(new-indv ``Clyde'' \{elephant\})} is functionally
equivalent to \textbf{(new-indv \{Clyde\} \{elephant\})}.

\subsubsection{The :english Keyword
Argument}\label{the-english-keyword-argument}

Most of the new-functions take an \textbf{:english} keyword argument.
This takes a list of names and syntax tags that is interpreted as in the
\textbf{english} function. So, for example, we might see a form like
this:

\textbf{(new-type \{swamp\} \{geographical area\}\\
:english \textquotesingle(:iname "wetland" "marsh" :adj "soggy"
"marshy"))}

This creates new type of land area with internal name \{swamp\} and
English names "swamp", "wetland", "marsh", and adjective names "soggy"
and "marshy".

By default, the new-element\textquotesingle s internal name becomes its
English name (or one of them), with a syntax-tag that is the default for
elements of this type: \textbf{:noun} for most indv and type nodes,
\textbf{:role} for new roles, and \textbf{:relation} for new relations.
This is in addition to whatever names are given by the \textbf{:english}
keyword argument. However, if the first element of the \textbf{:english}
argument is a syntax-tag, the iname gets that syntax tag instead of the
default. If you don\textquotesingle t want the internal name to become
an English name at all, put the tag \textbf{:no-iname} in the
\textbf{:english} list. As a rule, internal names that the Scone system
constructs itself are not saved as English names.

So, for example, if we execute the form

\textbf{(add-indv \{gizmo\} \{frobnitz\} :english ("foo" "bar"))}

the new \textbf{\{gizmo\}} node would have English names
\textbf{"gizmo"}, \textbf{"foo"} and \textbf{"bar"}, all with the
syntax-tag \textbf{:noun}. But for the form

\textbf{(add-indv \{noisy\} \{sound level\} :english
\textquotesingle(:adj "loud"))}

the new \textbf{\{noisy\}} node has English names \textbf{"noisy"} and
\textbf{"loud"}, both with syntax-tag \textbf{:adj}.

\subsubsection{Negation Links}\label{negation-links}

For each of the link-types in the system (except for the cancel-link),
there is a negated version: \textbf{is-not-a}, \textbf{not-eq}, and so
on. There is also a negated form of the generic statement link:
\textbf{not-statement}. If there is a negated link from A to B, that
will have the effect of over-riding any positive link of the
corresponding type from A to B. That is, the negated relation is always
stronger than the positive form of the relation whether the positive
form is inherited or stated directly. If you want to re-assert the
positive relation between A and B, you must first cancel the negative
link using a \textbf{cancel-link}.

Cancellation and negation in Scone is another topic that is beyond the
scope of this manual. It interacts in complex ways with transitive
relations. A quick rule of thumb is that a non-transitive statement can
either be cancelled or over-ridden by a negation; a transitive statement
can only be over-ridden, since there may be multiple paths; a negation
can only be cancelled (which is allowed, because negation links are
never transitive).

\subsection{Indv Nodes}\label{indv-nodes}

An \emph{individual node} (or \emph{indv-node}) represents an individual
entity rather than a type. It is not necessarily a physical object: it
can be something more abstract, such as a number or a specific date.
Individuals are normally the leaf-nodes of Scone\textquotesingle s is-a
hierarchy.

An indv-node may be \emph{proper} or \emph{generic}. A proper individual
node represents a specific, first-class individual such as
\textbf{\{George Washington\}} or \textbf{\{Paraguay\}} or \textbf{\{21
March 1948\}}. Normally, proper individuals are assumed to be distinct
from one another, and it is an error to equate two of them. However, in
special cases we may choose to over-ride this prohibition, for example
to equate \textbf{\{Clark Kent\}} to \textbf{\{Superman\}}.

A generic individual node represents some role or defined individual,
such as "Clyde\textquotesingle s mother" or \textbf{"}the first elephant
on the moon\textbf{"}. Generic individuals can be given properties and
relations, just like proper individuals: we can say things like
"Clyde\textquotesingle s mother is beautiful" or "the first elephant on
the moon was an African elephant". However, generic individuals are not
assumed to be distinct from one another -- in fact, we may think of a
generic individual as a stand-in for an proper individual to be named
later. It is legal to equate \textbf{\{mother of Clyde\}} to the proper
indv-node \textbf{\{Bertha T. Elephant\}} or to another generic
indv-node like \textbf{\{the first elephant on the moon\}} or to both of
these at once.

Some individual nodes are \emph{primitive individual nodes}. These Scone
elements represent Common Lisp objects such as numbers (integers,
ratios, or real numbers), strings, or functions. There are new-functions
for the various types of primitive nodes, but you normally create them
at read-time by putting the desired Lisp object within curly brackets:
\textbf{\{3\}}, \textbf{\{22/7\}}, \textbf{\{3.14159\}},
\textbf{\{"xyzzy"\}}, or \textbf{\{\#\textquotesingle append\}}. These
primitive nodes have a default parent based on their type -- a Scone
type-node such as \textbf{\{string\}} or \textbf{\{ratio\}}, but the
user may optionally provide a more specific parent node via the
\textbf{:parent} argument.

\subsubsection{Functions}\label{functions-1}

\textbf{new-indv} \emph{{[}FUNCTION{]}\\
}(iname parent \&key :context :proper :may-have :definition :english
:english-default)\\
This creates a new individual with the specified \textbf{iname} and
\textbf{parent}. (As noted above, if the \textbf{iname} argument is
\textbf{nil} the system will make up a unique name.)

The user may optionally specify a \textbf{:context} argument (defaults
to the value of \textbf{*context*}). By default, a has-link is also
created, connecting the context to the new indv-node.

If \textbf{:proper} is supplied and is explicitly \textbf{nil}, this is
a generic individual; otherwise assume the new node represents a proper
individual.

The \textbf{:may-have} argument, if T, indicates that the element may or
may not exist in the specified context. In this case, no has-link is
created.

If a \textbf{:definition} is supplied, this is a defined indv node, and
the definition is saved. This is intended for internal use.

The \textbf{:english} argument is used to supply one or more English
names for the new element, as described above. The default syntax-tag
for these English names is normally \textbf{:noun}, but a different tag
can be supplied via the \textbf{:english-default} argument.

\textbf{new-number} (value \&key :parent) \emph{{[}FUNCTION{]}}\\
\textbf{new-integer} (value \&key :parent) \emph{{[}FUNCTION{]}}\\
\textbf{new-ratio} (value \&key :parent) \emph{{[}FUNCTION{]}}\\
\textbf{new-float} (value \&key :parent) \emph{{[}FUNCTION{]}}\\
\textbf{new-string} (value \&key :parent) \emph{{[}FUNCTION{]}}\\
\textbf{new-function} (value \&key :parent) \emph{{[}FUNCTION{]}}

\textbf{new-struct} (value \&key :parent) \emph{{[}FUNCTION{]}}\\
Each of these functions create a new ``primitive'' individual node
representing some Common Lisp data object: a number (of any kind), an
integer, ratio, floating-point number, string, functions, or a structure
(also known as a ``defstruct''). The \textbf{value} argument must be a
Lisp object of the proper type. If another Scone object with this same
(eql) Lisp value already exists, we may return the existing object
instead of creating a new one. These objects are all created in the
\textbf{\{universal\}} context.

\textbf{node-value} (element) \emph{{[}FUNCTION{]}\\
}If \textbf{element} is a primitive individual node, extracts the value
from the iname and returns it as a Lisp object. Otherwise, returns
\textbf{nil}.

\subsection{Type Nodes}\label{type-nodes}

A \emph{type-node} represents the \emph{typical member} of some class or
set: \textbf{\{elephant\}}, \textbf{\{rainbow\}}, \textbf{\{weekday\}},
and so on. You can attach properties and class-memberships to a
type-node just as if it were an individual: "The typical elephant is a
mammal", "the color of the typical elephant is gray", and so on. In the
way they are used, type-nodes and indv-nodes are very similar, except
that individual nodes are normally leaf-nodes in the is-a hierarchy,
while type nodes are intended to have subtypes and instances.

Note that a type node represents the typical member of a set and not the
set itself. For every type node, there is an implicit \emph{set-node}
that does represent the set. However, we don\textquotesingle t actually
create the set node until/unless we have something to say about the set.
The \textbf{get-set-node} and \textbf{get-type-node} functions allow us
to access the set-node of a given type-node, and vice-versa. It is
perhaps confusing that the set-node of a type-node is an individual node
-- it represents an \emph{individual instance} of a set.

As with \textbf{new-indv}, the \textbf{new-type} function optionally
takes \textbf{:context} and \textbf{:english} arguments. The \textbf{:n}
argument is used to indicate the number of instances of this type that
exist in the specified (or default) context.

Most type-nodes represent description-based types such as ``elephant''.
For these, we have a description of the typical member, but not a
precise definition. However, Scone is also able to represent
\emph{defined-types}, which are normally created by the
\textbf{new-defined-type} function, described below. As new elements are
created and as the KB is modified, the Scone engine looks for members
and subtypes of each defined type. When one is found, the engine creates
a new \emph{derived is-a link}, placing the element under the defined
type in the is-a hierarchy. Once that is done, the new member may
inherit properties and relations that are true of the defined class, but
not part of the definition.\footnote{Because Scone's reasoning is not
  logically complete, we cannot guarantee that \emph{every} member of
  this defined class will be identified every the time the KB is
  modified. The computation involved may be intractable or unbounded.
  However, we will catch the vast majority of these. We plan to add
  additional inference mechanisms to the Scone engine to identify
  certain defined-class members at query time.}

\subsubsection{Functions}\label{functions-2}

\textbf{new-type} \emph{{[}FUNCTION{]}\\
}(iname parent \&key :context :definition :may-have :n :english
:english-default)\\
\textbf{\hfill\break
}Creates and returns a new type-node with the specified \textbf{iname}
(which will be created if the iname is \textbf{nil}) and
\textbf{parent}. By default, this also creates a has-link from the
specified context to the new type-node.

The user may optionally specify a \textbf{:context} argument (defaults
to the value of \textbf{*context*}). By default, a has-link is also
created, connecting the context to the new indv-node.

The \textbf{:may-have} argument, if T, indicates that the element may or
may not exist in the specified context. In this case, no has-link is
created.

If a has-link is created, the \textbf{:n} argument is placed on the
has-link to indicate how many elements of this type exist (or typically
exist) in the context.

If a \textbf{:definition} is supplied, this is a defined indv node, and
the definition is saved. This is intended for internal use.

The \textbf{:english} argument is used to supply one or more English
names for the new element, as described above. The default syntax-tag
for these English names is normally \textbf{:noun}, but a different tag
can be supplied via the \textbf{:english-default} argument.

\textbf{get-set-node} (e \&optional parent) \emph{{[}FUNCTION{]}\\
}Returns the set-node associated with type-node \textbf{e}, creating it
if necessary. The parent of the new set node is assumed to be
\textbf{\{set\}}, but an alternative value can be supplied as the
optional second argument.

\textbf{get-type-node} (e) \emph{{[}FUNCTION{]}\\
}Element \textbf{e} should be a set-node associated with some type-node.
Returns that type-node, if it exists, or \textbf{nil}.

\textbf{new-defined-type} (iname supertype-list predicate \&key :context
:english) \emph{{[}FUNCTION{]}}\\
Creates and returns a new defined-type node with the specified
\textbf{iname}. The definition is a predicate in two parts:
\textbf{supertype-list} is a list of elements that must all be superiors
of the candidate element. There must be at least one element in this
list, and for efficiency it should be as specific as possible. A single
object may be passed in instead of a list; \textbf{predicate}, if
non-nil, is a Lisp predicate function of one argument that must evaluate
to \textbf{t} for members of this defined class. The predicate should
not include testing the \textbf{supertype-list} -\/- that is done
separately.

\textbf{new-intersection-type} (iname parents \&key :context :english)
\emph{{[}FUNCTION{]}}\\
\textbf{parents} should be a list of elements, usually type-nodes.
Create and return a new defined-type element, with name \textbf{iname},
defined to be the intersection of all of the supertypes on the list.

\textbf{new-union-type} (iname parent subtypes \&key :context :english)
\emph{{[}FUNCTION{]}}\\
\textbf{subtypes} is a list of type-nodes or individuals. Creates and
returns a new type that is, by definition, the union of these subtypes.
All of these subtypes become members of a complete-split, indicating
that a member of the union-type must be a member of one -- and only one
-- of the subtypes.

\subsection{Map Nodes}\label{map-nodes}

A map-node represents, by definition, ``the x of y'' -- for example, the
\textbf{\{mother\}} of \textbf{\{Clyde\}}. Users do not normally deal
with map-nodes directly -- they are created and accessed by the
functions in section 7.13 below.

If the \textbf{role} argument is an individual role, the map-node is an
individual map node; if the role is a type-role, the map-node is a type
map node.

\subsubsection{Functions}\label{functions-3}

\textbf{new-map} (role owner \&key :iname :english :no-supplement)
\emph{{[}FUNCTION{]}}\\
Make and return a new map-node representing the \textbf{role} of the
\textbf{owner}. The role and owner must already exist -- no deferrals
are allowed here. Optionally provide an \textbf{:iname} and
\textbf{:english} names. If no \textbf{:iname} is supplied, we make up a
suitable one. The :no-supplement arg, if T, turns off some internal
housekeeping that is normally done.

\subsection{IS-A Links}\label{is-a-links}

An is-a link says that element A "is a" B in the specified context, and
that A should (by default) inherit all the properties and
type-memberships of B. We speak of B as being "above" A in the is-a
hierarchy.

If A and B are both type-nodes, this is a subtype relationship; if A is
an individual and B is a type, it is an instance-of relationship. The
type hierarchy branches in both directions (multiple inheritance), so a
node may have any number of incoming and outgoing is-a links. As noted
earlier, a node\textquotesingle s parent wire may take the place of one
of these outgoing is-a links.

\subsubsection{Functions}\label{functions-4}

\textbf{new-is-a} (a b \&key :negate :dummy :derived :iname :parent
:context :english) \emph{{[}FUNCTION{]}}\\
Creates and returns an is-a link from \textbf{a} to \textbf{b}.

\textbf{new-is-not-a} (a b \&key :iname :dummy :parent :context
:english) \emph{{[}FUNCTION{]}}\\
Make and return a new is-not-a link with the specified elements on the A
and B wires.

\subsection{EQ Links}\label{eq-links}

An \textbf{eq} link says that A and B are equivalent -- they represent
the same entity -- in the specified context. A inherits all the
properties of B and B inherits all the properties of A. It is common to
equate a proper individual with a generic individual, or two generic
individuals to one another. The system will complain if you try to
equate two proper individuals, though you can ignore the warning and
force this.

\subsubsection{Functions}\label{functions-5}

\textbf{new-eq} (a b \&key :negate :dummy :iname :parent :context
:english :no-supplement) \emph{{[}FUNCTION{]}}\\
Creates and returns an eq link from \textbf{a} to \textbf{b}.

\textbf{new-not-eq} (a b \&key :iname :dummy :parent :context :english)
\emph{{[}FUNCTION{]}}\\
Make and return a new not-eq link with the specified elements on the
\textbf{a} and \textbf{b} wires.

\subsection{Has Links}\label{has-links}

Note: The details of the has-link and of quantification in general are
under active review, and probably will be changing very significantly in
future releases of Scone.

\subsubsection{Functions}\label{functions-6}

\textbf{new-has} (a b \&key :n :negate :dummy :iname :parent :context
:english) \emph{{[}FUNCTION{]}}\\
Make and return a new HAS-LINK with the specified elements on the A and
B wires. Optionally provide :PARENT and :CONTEXT. If :N is specified,
that is the cardinality, stored ont eh C-wire.

\textbf{new-has-no} (a b \&key :iname :dummy :parent :context :english)
\emph{{[}FUNCTION{]}}\\
Make and return a new HAS-NO-LINK with the specified elements on the A
and B wires. Optionally provide :PARENT and :CONTEXT. This is equivalent
to NEW-HAS with the :NEGATE flag.

\subsection{Cancel Links}\label{cancel-links}

Suppose the knowledge base has a \textbf{\{hates\}} statement from
\textbf{\{elephant\}} to \textbf{\{snake\}}, indicating that elephants
hate snakes -- or, more precisely, that the typical elephant hates the
typical snake. Suppose that Clyde is an elephant, but an atypical one:
he doesn\textquotesingle t hate snakes. We can represent this kind of
exception using a cancel link from \textbf{\{Clyde\}} to the
\textbf{\{hates\}} statement inherited from \textbf{\{elephant\}}. The
effect of the cancel link is to turn off the \textbf{\{hates\}}
statement -- to render it inoperative -- when Scone is considering the
properties of \textbf{\{Clyde\}}. The statement is still effective for
other elephants.

We could also cancel snake-hating for an entire sub-class of elephants,
such as \textbf{\{African elephant\}} by connecting the cancel-link from
that type-node to the hates statement. Or we could tie the cancel-link
to a context: perhaps in \textbf{\{snake fantasy world\}}, everything is
the same as in \textbf{\{general\}} -- the real-world context -- except
that elephants do not hate snakes.

Cancel links are intended to cancel non-transitive statements, but not
class-memberships. For that, use an is-not-a link.

\subsubsection{Functions}\label{functions-7}

\textbf{new-cancel} (a b \&key :iname :dummy :parent :context :english)
\emph{{[}FUNCTION{]}}\\
Creates and returns a cancel link from \textbf{a} to \textbf{b}.

\subsection{Relations}\label{relations}

Scone allows the user to create new \emph{relations} and then to
instantiate these using \emph{statement} links. A relation element is
similar in form to a link -- it has A, B, and (sometimes) C wires -- but
it doesn\textquotesingle t really make an assertion of any kind. It is
simply a definition or a template for statement links that do make an
assertion.

A relation has an internal name or \emph{iname} such as \textbf{\{taller
than\}} or \textbf{\{employs\}}, naming the relation that holds from A
to B.

The \textbf{new-relation} function creates a new relation-element whose
wires are connected to the specified \textbf{:a}, \textbf{:b}, and
(optionally) \textbf{:c} arguments. The \textbf{:a} element represents
the domain of the relation, or rather the typical member of the domain
set; the \textbf{:b} element represents the typical member of the range;
the \textbf{:c} element, if present, is some value that is a function of
a specific A and B. For example, a relation might represent the distance
between A and B, with the C role being the value.

A very common situation is that the caller of \textbf{new-relation} will
want to create a new individual with a specific parent to be used as the
A element of the new relation. The most convenient way to do this is to
supply an \textbf{:a-inst-of} keyword followed by the parent element,
instead of the \textbf{:a} keyword. That indicates that the new-relation
function should create a new individual role node with the specified
parent and attach this to the new relation's a-wire. The
\textbf{:b-inst-of} and :\textbf{c-inst-of} keywords perform the same
function for the relation's other wires.

So, for example the user might create a new relation as follows:

\textbf{(new-relation \{reports to\} :a-inst-of \{person\} :b-inst-of
\{person\})}

This says that when we create a \textbf{\{reports to\}} statement, both
the A and B elements must be instances of \textbf{\{person\}}. But not
every \textbf{\{person\}} is involved in a \textbf{\{reports to\}}
relation.

The parent wire of the relation element may connect to a more general
relation, forming a type hierarchy of relations. For example, the parent
of the \textbf{\{hates\}} relation might be \textbf{\{dislikes\}}, since
we think of "hate" as a more specific kind of "dislike" -- a
particularly intense dislike.

Relations may be marked as transitive and/or symmetric, and these
attributes are inherited by the relation\textquotesingle s subtypes and
instances (statements). If we have a symmetric relation such as
\textbf{\{touches\}}, and we know that A touches B, it is not necessary
to also state that B touches A. Scone deduces that. Similarly, if we
have a transitive relation, such as \textbf{\{taller than\}}, Scone will
deduce that if A is taller than B and B is taller than C, then A is
taller than C.

\subsubsection{Functions}\label{functions-8}

\textbf{new-relation} \emph{{[}FUNCTION{]}\\
} (iname \&key :a :a-inst-of :a-type-of :b :b-inst-of :b-type-of :c
:c-inst-of :c-type-of\\
:parent :symmetric :transitive :english :inverse)\\
Define a new relation element that can be instantiated by a statement
link. The \textbf{iname} is the name of the relation from \textbf{a} to
\textbf{b}. This relation may have another relation as a parent, or a
type-node. Normally the caller supplies elements as \textbf{:a},
\textbf{:b}, and \textbf{:c} arguments. However, for the common case
where you want to create a new generic node with a specified parent, you
can instead supply \textbf{:a-inst-of} or \textbf{:a-type-of}, and
similarly for \textbf{b} and \textbf{c}. Optionally provide a
\textbf{:parent} element for the relation itself. Optionally set the
\textbf{:symmetric} and \textbf{:transitive} properties of the relation.
(These properties refer to the relation from \textbf{a} to \textbf{b}.)
The \textbf{:english} argument, as usual, is the name of the relation
(or list of names) in the forward direction. The \textbf{:inverse}
argument is similar in form, but in the reverse direction.

\subsection{Statement Links}\label{statement-links}

Once a relation has been defined, we can create statement links that
instantiates that relation.

\subsubsection{Functions}\label{functions-9}

\textbf{new-statement} (a rel b \&key :c :context :negate :dummy :iname
:english) \emph{{[}FUNCTION{]}}\\
Creates and returns a statement link representing the statement "a rel
b". For example, we might say something like this:

\begin{quote}
\textbf{(new-statement \{Clyde\} \{taller than\} \{Ernie\})}

The \textbf{rel} argument must be a relation element or another
statement. The \textbf{a} and \textbf{b} arguments must be compatible
with the a and b elements of \textbf{rel}, respectively. If the superior
\textbf{rel} has a c-node, we can supply a \textbf{:c} argument to the
statement.
\end{quote}

\textbf{new-not-statement} (a rel b \&key :c :context :iname :dummy
:english) \emph{{[}FUNCTION{]}}\\
Make and return a new not-statement with the specified elements on the
\textbf{a} and \textbf{b} wires.

\subsection{Split and Complete-Split
Links}\label{split-and-complete-split-links}

A \emph{split} is a special kind of link: it can have any number of
\emph{split-wires} instead of individual A, B, and C wires. The meaning
is that the type nodes connected by the split wires are mutually
disjoint: no individual can be an instance of more than one of these
types. If some of the connected nodes are individuals, they are disjoint
from one another and are not instances of any types in the split. Scone
has efficient machinery for determining whether a new or proposed link
causes a split to be violated.

So, for example, we might say

\begin{quote}
\textbf{(new-split \textquotesingle(\{bird\} \{reptile\} \{mammal\}))}
\end{quote}

to indicate that these are disjoint types with no common members.
However, there is no claim that the split types exhaust all the
possibilities. There may be other subtypes of \textbf{\{vertebrate\}}
that are not listed here.

A \emph{complete split} is like a split, but it makes the additional
claim that the split classes, taken together, cover all the members of
some supertype. Put another way, every member of the supertype must
belong to exactly one of the types in the complete split. So we might
say something like this:

\begin{quote}
\textbf{(new-complete-split \{vertebrate\}\\
\textquotesingle(\{fish\} \{amphibian\} \{reptile\} \{bird\}
\{mammal\})}
\end{quote}

A given type may have multiple splits and complete splits under it:

\begin{quote}
\textbf{(new-complete-split \{person\} \textquotesingle(\{male\}
\{female\}))\\
(new-complete-split \{person\} \textquotesingle(\{child\} \{adolescent\}
\{adult\}))}
\end{quote}

Once a split is in place, and attempt to violate the split will be
detected. However, it is possible to use a cancel-link to turn off the
split for a particular individual if, for example, we want to create an
individual who is both a male and a female. (However, if we do this, we
may have to use a number of additional cancellations to deal with
conflicting properties of this hybrid entity.)

\subsubsection{Functions}\label{functions-10}

\textbf{new-split} (members \&key :iname :dummy :parent :context
:english) \emph{{[}FUNCTION{]}}\\
The \textbf{members} argument is a list containing any number of
elements, usually type and individual nodes. Create and return a split
element stating that all of the members are distinct and disjoint.

\textbf{new-complete-split} \emph{{[}FUNCTION{]}\\
}(supertype members \&key :iname :dummy :parent :context :english)\\
Like new-split, but also states that any instance of the
\textbf{supertype} must belong to exactly one of the \textbf{members}.

\textbf{add-to-split} (split new-item) \emph{{[}FUNCTION{]}}\\
Add one \textbf{new-item} to the members of the specified
\textbf{split}, which may be a complete split.

\subsection{Functions to Create Multiple
Nodes}\label{functions-to-create-multiple-nodes}

These functions create multiple new subtypes or members in a given
class.

In each case, the \textbf{members} argument is a list of items. An item
may be an internal name for one of the new elements in the group, or a
list whose first element is the internal name; in that case, the
remainder of the list is treated as an \textbf{:english} argument for
this new element. By default, the iname becomes one of the English names
for the element, with syntax-tag \textbf{:noun}. If the first element of
the English list is a syntax-tag, the iname gets that tag instead. If
you don\textquotesingle t want the iname to be used as an English name
at all, include \textbf{:no-iname} in the list.

\subsubsection{Functions}\label{functions-11}

\textbf{new-split-subtypes} (parent members \&key :iname :context
:english) \emph{{[}FUNCTION{]}}\\
Create a set of disjoint subtypes under \textbf{parent}. First we create
a new split under \textbf{parent}. The \textbf{members} argument is a
list of internal names, one for each new type-node that is to be
created, as described above. Optionally, provide the \textbf{:context},
\textbf{:iname}, and \textbf{:english} names for the new split element.
Returns the new split element.

\textbf{new-complete-split-subtypes} (parent members \&key :iname
:context :english) \emph{{[}FUNCTION{]}}\\
Like \textbf{new-split-subtypes}, but creates a complete split of the
elements under \textbf{parent}.

\textbf{new-members} (parent members \&key :iname :context :english)
\emph{{[}FUNCTION{]}}\\
Like \textbf{new-split-subtypes}, but the items become new proper
individual nodes rather than subtypes of \textbf{parent}. Creates and
returns a split element stating that all these new indv-nodes are
distinct.

\textbf{new-complete-members} (parent members \&key :iname :context
:english) \emph{{[}FUNCTION{]}}\\
Like \textbf{new-complete-split-subtypes}, but the items become new indv
nodes rather than subtypes of \textbf{parent}. In other words, this
states that the items list contains all the individual members of the
parent set. Returns the new split element.

\subsection{Roles}\label{roles}

Roles are associated with type-nodes by the \textbf{new-indv-role} and
\textbf{new-type-role} functions. Once a role has been created for some
type, other functions described in this section can be used to attach
roles and values to instances of that type.

Note: Quantification machinery (\textbf{:n} and \textbf{:may-have}
arguments) for \textbf{new-indv-role} and \textbf{new-type-role} is
messed up, and will soon be replaced. For now, this correctly describes
what the code does.

\textbf{new-indv-role} \emph{{[}FUNCTION{]}\\
}(iname owner parent \&key :may-have :transitive :symmetric :english
:inverse)\\
Creates a new generic indv-node representing some role of the
\textbf{owner}. Also creates a \textbf{has} statement indicating that
for every instance of \textbf{owner}, there exists one instance of this
role. The \textbf{iname} is the name of the role-node. The
\textbf{parent} argument is the parent of the new role-node. The
\textbf{:english} argument is used to register English names for the
role, with the default syntax-type of \textbf{:role}. By default, the
iname becomes one of the English names.

\begin{quote}
If \textbf{:may-have} is present and non-null, that indicates that not
every instance of \textbf{owner} has this \textbf{role}. In this case,
we do not create the has-link.

The \textbf{:symmetric} and \textbf{:transitive} flags refer to the
implicit relation created between the owner and the role-node. The
\textbf{:english} argument, as usual, is the name of the role (or set of
names) in the forward direction. The \textbf{:inverse} argument is
similar in form, but in the reverse direction.
\end{quote}

\textbf{new-type-role} \emph{{[}FUNCTION{]}\\
}(iname owner parent \&key :n :may-have :transitive :symmetric :english
:inverse)\\
Like \textbf{new-indv-role}, but this function creates a type-role. The
value of \textbf{:n} is the number of instances of this role that the
owner typically has. If \textbf{:may-have} is present and non-null, it
indicates that not every instance of owner has this role, so we do not
create the has-link. Only one of \textbf{:n} and \textbf{:may-have}
should be supplied.

\textbf{x-is-the-y-of-z} (x y z) \emph{{[}FUNCTION{]}}\\
The function \textbf{x-is-the-y-of-z} finds or creates the node
representing ``the y of z''. Check whether \textbf{x} can be equated to
this node without causing any errors. If so, create an eq-link between
the role node and \textbf{x}. Return the new eq-link.

\textbf{x-is-a-y-of-z} (x y z) \emph{{[}FUNCTION{]}}\\
The function \textbf{x-is-a-y-of-z} finds or creates the node
representing "the y of z". This should be a type-node. Check whether
\textbf{x} can be of this type. If so, create and return an is-a link
from \textbf{x} to this node.

\textbf{the-x-of-y-is-a-z} (x y z) \emph{{[}FUNCTION{]}}\\
The function \textbf{the-x-of-y-is-a-z} finds or creates the node
representing "the x of y". Check whether this node can be of type
\textbf{z}. If so, create an is-a link from this node to type-node
\textbf{z}. Return the new is-a link.

\textbf{x-is-not-the-y-of-z} (x y z) \emph{{[}FUNCTION{]}}

\textbf{x-is-not-a-y-of-z} (x y z) \emph{{[}FUNCTION{]}}

\textbf{the-x-of-y-is-not-a-z} (x y z) \emph{{[}FUNCTION{]}}\\
These functions are like their positive counterparts, but they create a
not-eq or not-is-a link that will over-ride any existing relationship
and will prevent a new one from being formed (unless this not-link is
later cancelled).

\section{Removing Elements from the Scone
KB.}\label{removing-elements-from-the-scone-kb.}

\textbf{Beware:} Wholesale and random removal of elements can leave the
knowledge base in a bad state, with lots of dangling connections that
could give rise to anomalous behavior. However, it is sometimes
necessary to remove a few elements to repair some error or problem in
the knowledge base. It is usually safe to remove the last N elements
created, since these refer to earlier elements but earlier elements
usually don't refer to them, except in rare cases where the
deferred-connection mechanism has been used.

\subsubsection{Functions and Variables}\label{functions-and-variables-6}

\textbf{remove-element} (e) \emph{{[}FUNCTION{]}}\\
Disconnect element \textbf{e} from the network and remove it.

\textbf{remove-last-element} nil \emph{{[}FUNCTION{]}}\\
Remove the last element created.

\textbf{remove-elements-after} (e) \emph{{[}FUNCTION{]}}\\
Given some element \textbf{e}, remove all elements added later than
\textbf{e}, starting with the last one.

\textbf{remove-last-loaded-file} nil \emph{{[}FUNCTION{]}}\\
Remove all the elements, starting with the last-created, until we have
removed all elements created by the last loaded file (plus any created
interactively after that). Pops the last file off *LOADED-FILES* and
*LAST-LOADED-ELEMENTS*. Returns the number of elements remaining in the
KB.

\textbf{remove-currently-loading-file} nil \emph{{[}FUNCTION{]}}\\
Remove all the elements, starting with the last-created, until we have
removed all elements created by the file that is currently being loaded.
This is useful for backing out of a file-load if an error aborts the
load. Returns the number of elements remaining in the KB.

\textbf{remove-context} (c) \emph{{[}FUNCTION{]}}\\
The \textbf{c} argument is an element representing a context. Remove
\textbf{c}, all its sub-contexts, and all their contents. Return the
number of elements that remain in the KB after this operation.

\section{Queries and Predicates}\label{queries-and-predicates}

\subsection{Element-Type Predicate
Functions}\label{element-type-predicate-functions}

These functions all take a single argument \textbf{e}. They return
\textbf{t} if e is an element of the specified type, and \textbf{nil}
otherwise. The indentation indicates a subtype relation. For example, if
\textbf{number-node?} is true of an element, \textbf{primitive-node?},
\textbf{indv-node?}, and \textbf{node?} will be true as well.

\textbf{node? (e)}

\textbf{indv-node? (e)}

\begin{quote}
\textbf{proper-indv-node? (e)}

\textbf{generic-indv-node? (e)}

\textbf{defined-indv-node? (e)}

\textbf{primitive-node? (e)}

\textbf{number-node? (e)}

\textbf{integer-node? (e)}

\textbf{ratio-node? (e)}

\textbf{float-node? (e)}

\textbf{string-node? (e)}

\textbf{function-node? (e)}

\textbf{struct-node? (e)}
\end{quote}

\textbf{type-node? (e)}

\textbf{defined-type-node?}

\textbf{map-node? (e)}

\begin{quote}
\textbf{indv-map-node? (e)}

\textbf{type-map-node? (e)}
\end{quote}

\textbf{role-node? (e)}

\begin{quote}
\textbf{indv-role-node? (e)}

\textbf{type-role-node? (e)\\
}
\end{quote}

\textbf{link? (e)}

\textbf{is-a-link? (e)}

\textbf{is-not-a-link? (e)}

\textbf{eq-link? (e)}

\textbf{not-eq-link? (e)}

\textbf{has-link? (e)}

\textbf{has-no-link? (e)}

\textbf{cancel-link? (e)}

\textbf{statement? (e)}

\textbf{not-statement? (e)}

\textbf{split? (e)}

\begin{quote}
\textbf{complete-split? (e)}
\end{quote}

\textbf{relation? (e)}

The following predicates do not test the basic type of the element
\textbf{e}, but test some other feature of the element.

\textbf{defined?} (e) \emph{{[}FUNCTION{]}}\\
This predicate tests whether \textbf{e} is an element (type-node or
indv-node) with a specific definition, rather than just a description.

\textbf{derived?} (e) \emph{{[}FUNCTION{]}}\\
This predicate tests whether \textbf{e} is a derived element -- that is,
one created by the Scone engine rather than specifically by the user.

\textbf{killed?} (e) \emph{{[}FUNCTION{]}}\textbf{\hfill\break
}The killed flag on element \textbf{e} indicates that it is present but
inoperative.

\textbf{symmetric?} (e) \emph{{[}FUNCTION{]}}\textbf{\hfill\break
}The symmetric flag on a relation or statement indicates that the
relation holds in both directions.

\textbf{transitive?} (e) \emph{{[}FUNCTION{]}}\textbf{\hfill\break
}The transitive flag on a role or statement indicates that this is a
transitive relation.

\subsection{Predicates on the Is-A
Hierarchy}\label{predicates-on-the-is-a-hierarchy}

These predicates test for subtype, instance, and equality relations in
the is-a hierarchy.

\subsubsection{Functions and Variables}\label{functions-and-variables-7}

\textbf{is-x-a-y?} (x y) \emph{{[}FUNCTION{]}}\\
The \textbf{x} and \textbf{y} arguments are elements. This is a
predicate to determine whether \textbf{x} is known to be an instance or
subtype of \textbf{y} in the current context. A return of \textbf{:yes}
means \textbf{x} is definitely a \textbf{y}. \textbf{:no} means that
\textbf{x} is definitely not a \textbf{y}, and stating otherwise would
cause an error. \textbf{:maybe} means that we have no definite knowledge
either way.\\
\strut \\
A \textbf{:no} response will return an element as second value. This is
the particular split or superior type node that produces a conflict, or
one of these elements if there are multiple conflicts.

\textbf{simple-is-x-a-y?} (x y) \emph{{[}FUNCTION{]}}\\
The \textbf{x} and \textbf{y} arguments are elements. This is a
predicate to determine whether \textbf{x} is definitely known to be an
instance or subtype of \textbf{y} in the current context. Returns T or
NIL.

\textbf{is-x-eq-y?} (x y) \emph{{[}FUNCTION{]}}\\
The \textbf{x} and \textbf{y} arguments are elements. This is a
predicate to determine whether \textbf{x} is known to be equivalent to
\textbf{y} in the current context. A return of \textbf{:yes} means
\textbf{x} is definitely equivalent to \textbf{y}. \textbf{:no} means
that \textbf{x} is definitely not equivalent to \textbf{y}, and stating
otherwise would cause an error. \textbf{:maybe} means that we have no
definite knowledge either way.\\
\strut \\
A \textbf{:no} response will return an element as second value. This is
the particular split or element node that produces a conflict, or one of
these elements if there are multiple conflicts.

\textbf{simple-is-x-eq-y?} (x y) \emph{{[}FUNCTION{]}}\\
The \textbf{x} and \textbf{y} arguments are elements. This is a
predicate to determine whether \textbf{x} is definitely known to be
equivalent to \textbf{y} in the current context. Returns T or NIL.

\subsection{Mark, List, and Show
Functions}\label{mark-list-and-show-functions}

Most of these query functions have three versions: a \emph{mark}
version, which marks all the members of some class (such as "instances
of mammal") with a specified marker; a \emph{list} version", which
returns a list of such elements, suitable for processing by an external
Lisp function; and a \emph{show} version that prints the set of elements
in a human-friendly format.

The "show" functions all take some additional keyword arguments that
control the printing:

The \textbf{:n} argument, an integer, is the maximum number of items
that will be printed. The default is 100.

An \textbf{:all} argument of \textbf{t} over-rides the \textbf{:n}
argument and forces printing of all the elements in the result set. The
default is \textbf{nil}.

The \textbf{:last} argument, an integer, says to print the last
\textbf{n} arguments in the set rather than the first \textbf{n}.

Normally you call a show function with only one of \textbf{:n},
\textbf{:all}, or \textbf{:last}.

\subsubsection{Functions and Variables}\label{functions-and-variables-8}

\textbf{list-elements} nil \emph{{[}FUNCTION{]}}

\textbf{show-elements} (\&key :n :all :last) \emph{{[}FUNCTION{]}}\\
List or print all the elements in the KB, in the order in which the
elements were defined.

\textbf{list-marked} (m) \emph{{[}FUNCTION{]}}

\textbf{show-marked} (m \&key :n :all :last :heading)
\emph{{[}FUNCTION{]}}\\
List or print all the elements marked with marker \textbf{m}, in the
order in which the markers were placed.

\textbf{list-not-marked} (m) \emph{{[}FUNCTION{]}}

\textbf{show-not-marked} (m \&key :n :all :last) \emph{{[}FUNCTION{]}}\\
List or print all the elements not marked with marker \textbf{m}.

\textbf{mark-supertypes} (e m) \emph{{[}FUNCTION{]}}\\
\textbf{list-supertypes} (e) \emph{{[}FUNCTION{]}}

\textbf{show-supertypes} (e \&key :n :all :last) \emph{{[}FUNCTION{]}\\
}Mark/list/print all the types above \textbf{e} in the is-a hierarchy
(in the current context).

\textbf{mark-subtypes} (e m) \emph{{[}FUNCTION{]}\\
}\textbf{list-subtypes} (e) \emph{{[}FUNCTION{]}}

\textbf{show-subtypes} (e \&key :n :all :last) \emph{{[}FUNCTION{]}}\\
Mark/list/print all the types below \textbf{e} in the is-a hierarchy (in
the current context).

\textbf{mark-instances} (e m) \emph{{[}FUNCTION{]}}\\
\textbf{list-instances} (e) \emph{{[}FUNCTION{]}}

\textbf{show-instances} (e \&key :n :all :last) \emph{{[}FUNCTION{]}}\\
Mark/list/print all the instances (indv nodes) under \textbf{e} in the
current context.

\textbf{mark-parents} (e m) \emph{{[}FUNCTION{]}}\\
\textbf{list-parents} (e) \emph{{[}FUNCTION{]}}

\textbf{show-parents} (e \&key :n :all :last) \emph{{[}FUNCTION{]}}\\
Mark/list/print the immediate superiors of element \textbf{e} in the
is-a hierarchy (in the current context).

\textbf{mark-children} (e m) \emph{{[}FUNCTION{]}}\\
\textbf{list-children} (e) \emph{{[}FUNCTION{]}}

\textbf{show-children} (e \&key :n :all :last) \emph{{[}FUNCTION{]}}\\
Mark/list/print the immediate descendants of element \textbf{e} in the
is-a hierarchy (in the current context).

\textbf{mark-superiors} (e m) \emph{{[}FUNCTION{]}}\\
\textbf{list-superiors} (e) \emph{{[}FUNCTION{]}}

\textbf{show-superiors} (e \&key :n :all :last) \emph{{[}FUNCTION{]}\\
}Mark/list/print all the elements, regardless of type, above \textbf{e}
in the is-a hierarchy (in the current context).

\textbf{mark-inferiors} (e m) \emph{{[}FUNCTION{]}}\\
\textbf{list-inferiors} (e) \emph{{[}FUNCTION{]}}

\textbf{show-inferiors} (e \&key :n :all :last) \emph{{[}FUNCTION{]}}\\
Mark/list/print all the elements, regardless of type, below \textbf{e}
in the is-a hierarchy (in the current context).

\textbf{mark-equals} (e m) \emph{{[}FUNCTION{]}}\\
\textbf{list-equals} (e) \emph{{[}FUNCTION{]}}

\textbf{show-equals} (e \&key :n :all :last) \emph{{[}FUNCTION{]}}\\
Mark/list/print all the elements, regardless of type, that are
equivalent to \textbf{e} (i.e. connected by an EQ-link or chain of them)
in the is-a hierarchy (in the current context).

\textbf{mark-intersection} (superiors m) \emph{{[}FUNCTION{]}}\\
\textbf{list-intersection} (superiors) \emph{{[}FUNCTION{]}}

\textbf{show-intersection} (superiors \&key :n :all :last)
\emph{{[}FUNCTION{]}\\
}Given list of \textbf{superiors}, which should all be type-nodes,
mark/list/print every element that is in the intersection of all these
superior types. Return the count.

\textbf{mark-lowermost} (m1 m2) \emph{{[}FUNCTION{]}\\
}\textbf{list-lowermost} (m) \emph{{[}FUNCTION{]}}

\textbf{show-lowermost} (m \&key :n :all :last) \emph{{[}FUNCTION{]}}\\
The \textbf{mark-lowermost} function examines the set of elements
currently marked with \textbf{m1} and places \textbf{m2} on the
lowermost elements of that set -- that is, on all elements marked with
\textbf{m1} that do not have an immediate descendant that is also marked
with \textbf{m1}. The \textbf{list-lowermost} and
\textbf{show-lowermost} functions list/print the lowermost elements in
the set marked with \textbf{m}.

\textbf{mark-uppermost} (m1 m2) \emph{{[}FUNCTION{]}}\\
\textbf{list-uppermost} (m) \emph{{[}FUNCTION{]}}

\textbf{show-uppermost} (m \&key :n :all :last) \emph{{[}FUNCTION{]}}\\
The \textbf{mark-uppermost} function examines the set of elements
currently marked with \textbf{m1} and places \textbf{m2} on the
uppermost elements of that set -- that is, on all elements marked with
\textbf{m1} that do not have an immediate superior that is also marked
with \textbf{m1}. The \textbf{list-uppermostmost} and
\textbf{show-uppermost} functions list/print the uppermost elements in
the set marked with \textbf{m}.

\textbf{mark-proper} (m1 m2) \emph{{[}FUNCTION{]}\\
}\textbf{list-proper} (m) \emph{{[}FUNCTION{]}}

\textbf{show-proper} (m \&key :n :all :last) \emph{{[}FUNCTION{]}}\\
From the set of elements currently marked with M (or M1),
mark/list/print any proper individual nodes that are lowermost in the
type hierarchy.

\textbf{mark-most-specific} (m1 m2) \emph{{[}FUNCTION{]}}\\
\textbf{list-most-specific} (m) \emph{{[}FUNCTION{]}}

\textbf{show-most-specific} (m \&key :n :all :last :heading)
\emph{{[}FUNCTION{]}}\\
The \textbf{mark-most-specific} function examines the set of elements
currently marked with \textbf{m1} and places \textbf{m2} on the most
specific elements of that set -- that is, any proper individual nodes
and, if none of those are present, the lowermost type or generic nodes..
The \textbf{list-most-specific} and \textbf{show-most-specific}
functions list/print the most specific elements in the set marked with
\textbf{m}.

\textbf{acceptable-role-filler} (m) \emph{{[}FUNCTION{]}}\\
From the set of nodes currently marked with M, find the one that would
serve best as an answer for THE-X-OF-Y and similar functions. Currently
we accept lowermost proper individuals, types, and statements. If more
than one exists, pick one at random; if none, return NIL.

\subsection{Query Functions for Roles and
Relations}\label{query-functions-for-roles-and-relations}

\textbf{mark-role} (role owner m \&key :downscan :augment)
\emph{{[}FUNCTION{]}}\\
OWNER is any element. ROLE is a node that is a role of OWNER, directly
or by inheritance. Place marker M on all nodes X such that
\textquotesingle X is the ROLE of OWNER\textquotesingle. If :AUGMENT,
don\textquotesingle t clear M before marking new nodes. Downscan the M
markers unless :DOWNSCAN is explicitly NIL. Returns the number of
elements ultimately marked with M.

\textbf{mark-role-inverse} (role player m \&key :downscan :augment)
\emph{{[}FUNCTION{]}}\\
Place marker M on all nodes X such that \textquotesingle PLAYER is the
ROLE of X\textquotesingle. If :AUGMENT, don\textquotesingle t clear M
before marking new nodes. Downscan the M markers unless :DOWNSCAN is
explicitly NIL. Returns the number of elements ultimately marked with M.

\textbf{mark-the-x-of-y} (x y m) \emph{{[}FUNCTION{]}}\\
Mark with M all elements E that collectively represent the X of Y. Do
not propagate M down to subtypes and instances of this role.

\textbf{mark-all-x-of-y} (x y m) \emph{{[}FUNCTION{]}\\
}\textbf{list-all-x-of-y} (x y) \emph{{[}FUNCTION{]}}

\textbf{show-all-x-of-y} (x y \&key :n :all :last)
\emph{{[}FUNCTION{]}}\\
Print out all the X of Y.

\textbf{the-x-of-y} (x y) \emph{{[}FUNCTION{]}}\\
Find the most specific proper node representing the X of Y, and return
that node if it exists. Else return NIL.

\textbf{mark-the-x-inverse-of-y} (x y m) \emph{{[}FUNCTION{]}}\\
Mark with M all elements E such that the X of E is Y. Do not propagate M
down to subtypes and instances of this role.

\textbf{mark-all-x-inverse-of-y} (x y m) \emph{{[}FUNCTION{]}}\\
\textbf{list-all-x-inverse-of-y} (x y) \emph{{[}FUNCTION{]}}

\textbf{show-all-x-inverse-of-y} (x y \&key :n :all :last)
\emph{{[}FUNCTION{]}}\\
Mark with \textbf{m}, list, or print out all nodes \textbf{e} such that
the \textbf{x} of \textbf{e} is \textbf{y}.

\textbf{the-x-inverse-of-y} (x y) \emph{{[}FUNCTION{]}}\\
Find the most specific proper node E such that the X of E is Y. If no
such E exists, return NIL. If there are more than one E, return one of
them.

\textbf{mark-roles} (e m) \emph{{[}FUNCTION{]}}\\
\textbf{list-roles} (e) \emph{{[}FUNCTION{]}}

\textbf{show-roles} (e \&key :n :all :last) \emph{{[}FUNCTION{]}}\\
Mark with \textbf{m}, list, or print out all the roles defined for
element \textbf{e} in the current context, whether or not those roles
are filled.

\textbf{the-x-role-of-y} (x y) \emph{{[}FUNCTION{]}}\\
Find and return the role-node or map-node directly representing the
X-role of Y. If this does not exist, create it.

\textbf{can-x-have-a-y?} (x y) \emph{{[}FUNCTION{]}}\\
Predicate to determine whether element X can have a Y role in the
current context. That is, return T if Y is a role of X or some node
above X, else NIL.

\textbf{can-x-be-the-y-of-z?} (x y z) \emph{{[}FUNCTION{]}}\\
Predicate to determine whether it would be legal to state that X is the
Y of Z. That is, would any splits or other constraints be violated? If
it is OK, return T. If it is not, return NIL with the unhappy split or
constraint as the second return value. If there are multiple problems,
we just return the first one we hit.

\textbf{can-x-be-a-y-of-z?} (x y z) \emph{{[}FUNCTION{]}}\\
Predicate to determine whether it would be legal to state that X is a Y
of Z. That is, would any splits or other constraints be violated? If it
is OK, return T. If it is not, return NIL with the unhappy split or
constraint as the second return value. If there are multiple problems,
we just return the first one we hit.

\textbf{statement-true?} (a rel b) \emph{{[}FUNCTION{]}\\
}A predicate that tests whether "\textbf{a rel b}" is true in the
current context. If so, the second return value is the link that
actually states this fact (if there is a single link that does this).
The third return value is that link's c-node, if any. If the predicate
is false, return only the single value \textbf{nil}.

\textbf{mark-rel} (rel a m \&key :downscan :augment
:recursion-allowance) \emph{{[}FUNCTION{]}}

\textbf{list-rel} (rel a) \emph{{[}FUNCTION{]}}

\textbf{show-rel} (rel a \&key :n :all :last) \emph{{[}FUNCTION{]}\\
}In \textbf{mark-rel}, the \textbf{a} argument is an element. The
\textbf{rel} argument is a relation-node. The \textbf{mark-rel} function
places marker \textbf{m} on all elements \textbf{b} for which the
statement "\textbf{a rel b}" is known to be true in the current context.
Then downscan marker \textbf{m} unless the \textbf{:downscan} argument
is explicitly \textbf{nil}. \textbf{:recursion-allowance} controls the
number of levels of descriptions we will explore. Returns the number of
elements marked with \textbf{m} after this operation. The
\textbf{list-rel} and \textbf{show-rel} functions list and print the set
of elements that \textbf{mark-rel} would mark.

\textbf{mark-rel-inverse} (rel a m \&key :downscan :augment
:recursion-allowance) \emph{{[}FUNCTION{]}\\
}\textbf{list-rel-inverse} (rel b) \emph{{[}FUNCTION{]}}

\textbf{show-rel-inverse} (rel b \&key :n :all :last)
\emph{{[}FUNCTION{]}}\\
Like \textbf{mark-rel} and friends, but here we cross the statement
links in the opposite direction. That is, we mark/list/print all
elements a such that "a rel b".

\subsection{Cardinality}\label{cardinality}

\emph{Cardinality} queries answer questions like ``does X have a Y?''
and ``how many Ys does X have?''. The engine source reserves a
\textbf{Cardinality} section (see \textbf{engine.lisp}) and sketches a
\textbf{does-x-have-a-y?} predicate built on \textbf{statement-true?},
but the code block is commented out and marked with a
\textbf{to-do} note: ``Implement the cardinality functions.'' As of
this release, there is no supported cardinality API in Scone.
Applications that need to count role fillers can do so by calling
\textbf{list-all-x-of-y} (or one of the similar listing functions in
the Queries section) and taking the length of the result.

\subsection{Miscellaneous Show
Functions}\label{miscellaneous-show-functions}

\textbf{show-element-counts} nil \emph{{[}FUNCTION{]}}\\
Print a table showing the number of elements of each type.

\textbf{show-marker-counts} nil \emph{{[}FUNCTION{]}}\\
Prints the number of elements marked with each marker.

\textbf{show-ontology} (\&key :under :n) \emph{{[}FUNCTION{]}}\\
Print every node in the ontology beneath \textbf{start}. This output is
designed to be read by humans, not programs. If the \textbf{:under}
argument (a Scone element) is supplied, print only the part of the
ontology under this element. If \textbf{:n} is supplied, stop after N
elements are printed.

\section{Marker Scans}\label{marker-scans}

The various marker scans are powerful and rather complex internal
functions that perform the pseudo-parallel search and inference in
Scone. Once Scone is mature, users will seldom have to deal directly
with these scans, but for now they provide the means for writing any
search and inference functions that Scone does not supply.

\subsubsection{Functions and Variables}\label{functions-and-variables-9}

\textbf{upscan} (start-element m \&key :mode :m1 :augment :restrict-m)
\emph{{[}FUNCTION{]}}\\
Mark \textbf{start-element} with marker \textbf{m}. Propagate this
marker upwards, across parent wires and is-a links that are active in
the current context, while respecting is-not-a links and cancellations.
Eq-links are crossed in both directions. Places the cancel-marker
corresponding to \textbf{m} on nodes that explicitly are not considered
to be superiors of start-element.

Returns the number of elements marked with \textbf{m} after the
\textbf{upscan}

\begin{quote}
Normally \textbf{upscan} clears the \textbf{m} marker and its
cancel-marker before start-element is marked and the scan is performed.
However, if the \textbf{:augment} argument is \textbf{t}, we leave any
pre-existing \textbf{m} markers in place and they are propagated upward
as well.

The \textbf{mode} argument is one of \textbf{:cross-map},
\textbf{:gate-map}, or \textbf{:no-map}. The default is
\textbf{:cross-map}. The others are specialized scan types with respect
to propagating across map connections. For some modes, we need an
additional marker, which is supplied as \textbf{:m1}.

The \textbf{:restrict-m} argument, if present, is a marker already
present on some elements in the network. In this scan we place
\textbf{m} only on elements that already have \textbf{restrict-m}.
Returns the number of elements marked with \textbf{m} after the scan.

The operation of upscan and the other scans may be modified by the
\textbf{*one-step*}, \textbf{*ignore-context*}, and
\textbf{*default-recursion-allowance*} variables.
\end{quote}

\textbf{*ignore-context*} \emph{{[}VARIABLE{]}}\\
If \textbf{t}, marker-propagations pay no attention to whether or not
the nodes and links are in an active context. This variable is used for
upscanning while activating a context, and may also be useful finding
whether some entity exists in \emph{any} context.

\textbf{*one-step*} \emph{{[}VARIABLE{]}}\\
If \textbf{t}, marker propagations go only one level in the desired
direction. Used in functions such as \textbf{mark-parents}, where you
don\textquotesingle t want the full transitive closure.

\textbf{*default-recursion-allowance*} \emph{{[}VARIABLE{]}}\\
In functions such as \textbf{mark-rel}, we must individually activate
and explore each description instance in which the source node plays a
role, and this process may recurse. This is the default allowance for
recursion depth that we will explore.

\textbf{downscan} (start-element m \&key :mode :m1 :augment :restrict-m)
\emph{{[}FUNCTION{]}}\\
Like upscan, but propagates marker \textbf{m} and its cancel-marker
downward from \textbf{start-element}. That is, the direction is from a
type to its subtypes and instances.

\textbf{eq-scan} (start-element m \&key :mode :m1 :augment :restrict-m)
\emph{{[}FUNCTION{]}}\\
Like upscan, but propagates marker \textbf{m} only to elements that are
equivalent to \textbf{start-element}. That is, this does not propagate
upward or downward in the type hierarchy, but only crosses eq links.

\textbf{lowermost?} (e m) \emph{{[}FUNCTION{]}}\\
Assume element \textbf{e} is marked with marker \textbf{m}. Determine
whether \textbf{e} is one of the lowermost nodes in this marked set --
that is, it has no immediate children also marked with \textbf{m}.

\textbf{uppermost?} (e m) \emph{{[}FUNCTION{]}}\\
Like \textbf{lowermost?}, but determines whether \textbf{e} is an
uppermost member of the set marked with \textbf{m}.

\textbf{splits-ok?} (m) \emph{{[}FUNCTION{]}}\\
Some set of elements has been marked with \textbf{m} and upscanned. Look
for splits that have \textbf{m} on more than one of the disjoint items.
If there are none, return \textbf{t}, meaning OK. If we find an unhappy
split, return \textbf{nil}, returning the unhappy split element (or one
of them, if there are more than one) as the second return value.

\textbf{which-split-member?} (e split) \emph{{[}FUNCTION{]}}\\
Given an element \textbf{e} and a split-node \textbf{split}, finds which
of the split type nodes \textbf{e} is an inferior of. If none of the
types is a supertype of \textbf{e}, returned \textbf{nil}.

\section{Miscellaneous Functions and
Variables}\label{miscellaneous-functions-and-variables}

\subsection{Commentary, Printing, and
Naming}\label{commentary-printing-and-naming}

\textbf{commentary} (\&rest args) \emph{{[}FUNCTION{]}}\\
This function is identical in form to the Common Lisp \textbf{format}
function, but without a stream argument. The \textbf{commentary}
function prints to \textbf{*commentary-stream*}, which may or may not be
equivalent to \textbf{*standard-output*}. If that variable is
\textbf{nil}, then \textbf{commentary} doesn\textquotesingle t print at
all.

\begin{quote}
The intent is that \textbf{commentary} should be used to inform the user
when the system has done something interesting, but with the capability
to turn off all such comments in non-interactive situations or if the
user doesn't want to see all this information.
\end{quote}

\textbf{*commentary-stream*} \emph{{[}VARIABLE{]}}\\
Stream to which \textbf{commentary} prints its output.

\textbf{*print-namespace-in-elements*} \emph{{[}VARIABLE{]}}\\
If \textbf{t}, always include the namespace, followed by a colon, in
element inames. This ensures that what we print can be read back in
without any namespace ambiguity. If \textbf{nil}, do not include the
namespace. This is better for human readabilty. If \textbf{:maybe},
print the namespace if it differs from the current \textbf{*namespace*}.
This is usually safe.

\textbf{*generate-long-element-names*} \emph{{[}VARIABLE{]}}\\
If \textbf{t}, generate long names for certain element types. For
example, a statement link will have an internal name that is an English
description of arguments and the relation. This is good for development,
but significantly increases the storage requirement for large knowledge
bases.

\textbf{*print-keylist*} \emph{{[}VARIABLE{]}}\\
If \textbf{t}, dumping an element pattern includes the list of
keyword/value pairs as part of the pattern. That is necessary if we want
to read that pattern back in, but we set this to \textbf{nil} for some
human-readable dumps.

\textbf{*show-ontology-indent*} \emph{{[}VARIABLE{]}}\\
An integer, the number of spaces to indent each level in
\textbf{show-ontology}.

\textbf{*show-stream*} \emph{{[}VARIABLE{]}}\\
Stream to which we print the output of \textbf{show} functions. If
\textbf{nil}, don\textquotesingle t print these at all.

\subsection{Variables Linking to Essential Scone
Elements}\label{variables-linking-to-essential-scone-elements}

There are a number of elements set up in the \textbf{bootstrap.lisp} KB
file that are referred to by code in the Scone engine. For convenience,
we store each of these elements in a global variable so that it is easy
for engine and user code to reference these elements. So, for example, a
pointer to the element \textbf{\{thing\}} is stored in the variable
\textbf{*thing*}.

\textbf{*thing*\\
}The uppermost node in the is-a hierarchy.

\textbf{*universal*\\
}The pre-defined \textbf{\{universal\}} context.

\textbf{*general*\\
}The pre-defined \textbf{\{general\}} context.

\textbf{*set*\\
*empty-set*}

\textbf{*non-empty-set*\\
}Pre-defined type nodes related to the concept of \textbf{\{set\}}.

\textbf{*cardinality*\\
}The \textbf{\{cardinality\}} role for the \textbf{\{set\}} node.

\textbf{*zero*\\
*one*}

\textbf{*two*\\
}Some pre-defined number nodes.

\textbf{*number*\\
*integer*\\
*ratio*\\
*float*\\
*string*\\
*function*\\
*struct*\\
*relation*\\
*is-a-link*\\
*is-not-a-link*\\
*eq-link*\\
*not-eq-link*\\
*cancel-link*\\
*has-link*\\
*has-no-link*\\
*split*\\
*complete-split*\\
*statement-link*}

\textbf{*not-statement-link*\\
}These type-nodes represent the various element-types built into Scone.

\section{Consistency Checking}\label{consistency-checking}

Scone does not in general try to prove that a KB is globally
consistent -- the open-world assumption and the ability to override
inherited information via cancel-links mean that ``inconsistency'' is
a local property. The engine does, however, supply a small set of
utilities for finding places where the is-a hierarchy conflicts with
declared \emph{splits} (disjoint sets) or \emph{cancel-links}. These
functions are useful for auditing a KB, for sanity-checking a newly
loaded file, and for deciding whether a proposed new link would
introduce a conflict before committing it.

All five functions below rely on the ordinary marker-scan machinery
and are safe to call at any point after the KB is loaded.

\subsubsection{Functions}\label{functions-consistency}

\textbf{find-split-violations} \emph{{[}FUNCTION{]}}\\
Scan every element in the KB. For each element, upscan to collect its
ancestors and test whether that ancestor set contains more than one
member of any declared split. For each uppermost offending element,
print the element together with the split it violates. Returns no
value; the report is printed to standard output. This is the function
to run after a large load to discover latent conflicts in the
ontology.

\textbf{violated-split?} (m) \emph{{[}FUNCTION{]}}\\
Given a marker \textbf{m} that has already been upscanned over some
starting element, return the first split-link whose disjoint members
are marked by more than one \textbf{m}-marked element, or \textbf{nil}
if no such split is found. This is the underlying primitive that
\textbf{find-split-violations} uses; applications that want to test a
single element for split conflicts can call it directly after an
\textbf{upscan}.

\textbf{split-violations-below?} (e1 e2) \emph{{[}FUNCTION{]}}\\
Test whether creating an IS-A or EQ link from \textbf{e1} to
\textbf{e2} would cause any element below \textbf{e1} in the is-a
hierarchy to violate a split. Any uppermost conflicted elements are
printed. Returns \textbf{t} if at least one violation is found,
\textbf{nil} otherwise. Use this to pre-check a proposed
\textbf{new-is-a} or \textbf{new-eq} before committing to it.

\textbf{incompatible?} (e1 e2) \emph{{[}FUNCTION{]}}\\
Test whether joining \textbf{e1} and \textbf{e2} with an IS-A or EQ
link would violate any split or cancel constraint. Upscans each
element with its own marker and looks for (a) any element that carries
one element's ancestry plus the cancel marker of the other element's
ancestry, and (b) any active split whose disjoint members straddle the
two ancestry sets. Returns the first offending element, or \textbf{nil}
if the join is clean. Unlike \textbf{split-violations-below?}, this
checks both directions and includes cancel-link conflicts as well as
splits.

\textbf{find-path} (lower upper) \emph{{[}FUNCTION{]}}\\
Print every element on an is-a path between \textbf{lower} and
\textbf{upper}. This is done by upscanning from \textbf{lower},
downscanning from \textbf{upper}, and printing the intersection. Not a
consistency check in itself, but invaluable for diagnosing why a
predicate is returning an unexpected answer.

\section{Structure-Creation Functions from the Core
KB}\label{structure-creation-functions-from-the-core-kb}

These commonly-used functions for creating new network structure are not
defined in the Scone engine file, but instead are defined in the files
contained in the ``core'' knowledge base. They must be defined as part
of the knowledge base because they build upon element structures created
as the KB files are being loaded.

\subsubsection{Functions and
Variables}\label{functions-and-variables-10}

\textbf{new-measure} (magnitude unit) {[}FUNCTION{]}\\
This function, defined in the ``units'' file, creates and returns an
individual of type \textbf{\{measure\}}, which is the combination of a
\textbf{magnitude} (i.e. a number) and a \textbf{unit}, which should be
an instance of \textbf{\{unit\}}. So the measure represents something
like ``16.0 ton''. The \textbf{magnitude} may be an integer, ratio, or a
float, or a Scone object of type \textbf{\{number\}}. A specific
sub-type of \textbf{\{measure\}} is chosen based on the type of
\textbf{unit}.

\begin{quote}
Note that the \textbf{unit} should be a singular form: \textbf{"ton"}
rather than \textbf{``tons''}. Scone does not yet have the machinery to
deal smoothly with plurals and other natural-language morphology.
\end{quote}

\textbf{new-quantity} (magnitude unit stuff) {[}FUNCTION{]}\\
This is like \textbf{new-measure}, but takes a third argument
\textbf{stuff}, which should be a subtype of the type
\textbf{\{stuff\}}, which is defined in the ``units'' file. It creates
and returns an element representing something like ``16.0 ton of coal''
or ``10 megabytes of census-data''. Note that \textbf{\{stuff\}}
represents any quality measurable by a unit, and not necessarily a
physical substance, which in Scone is a \textbf{\{material\}}.

\section{Alphabetical Index of Functions and
Variables}\label{alphabetical-index-of-functions-and-variables}

*activation-cancel-marker* 27

*activation-marker* 27

*cancel-link* 51

*cardinality* 51

*check-defined-types* 9

*commentary-stream* 50

*comment-on-defined-types* 9

*comment-on-element-creation* 9

*complete-split* 51

*context* 27

*context-cancel-marker* 27

*context-marker* 27

*create-undefined-elements* 9

*currently-loading-files* 8

*deduce-owner-type-from-role* 9

*default-kb-pathname* 8

*default-recursion-allowance* 48

*deferred-connections* 10

*defer-unknown-connections* 10

*disambiguate-policy* 9

*empty-set* 51

*eq-link* 51

*first-element* 14

*float* 51

*function* 51

*general* 51

*generate-long-element-names* 50

*has-link* 51

*has-no-link* 51

*ignore-context* 48

*integer* 51

*is-a-link* 51

*is-not-a-link* 51

*kb-logging-stream* 12

*last-element* 14

*legal-syntax-tags* 21

*loaded-files* 8

*loading-kb-file* 8

*namespace* 19

*namespaces* 19

*n-available-markers* 26

*n-elements* 14

*no-kb-error-checking* 9

*non-empty-set* 51

*not-eq-link* 51

*not-statement-link* 51

*number* 51

*one* 51

*one-step* 48

*print-keylist* 50

*print-namespace-in-elements* 50

*ratio* 51

*relation* 51

*reloading-kb-file* 9

*set* 51

*show-ontology-indent* 50

*show-stream* 50

*split* 51

*statement-link* 51

*string* 51

*struct* 51

*thing* 51

*two* 51

*universal* 51

*verbose-loading* 8

*zero* 51

acceptable-role-filler 45

add-to-split 37

a-element 15

b-element 15

cancel-link? 41

can-x-be-a-y-of-z? 46

can-x-be-the-y-of-z? 46

can-x-have-a-y? 46

c-element 15

check-legal-marker 24

check-legal-marker-pair 24

check-loaded-files 11

checkpoint-kb 11

checkpoint-new 11

clear-all-markers 25

clear-element-property 16

clear-marker 25

clear-marker-pair 25

commentary 50

complete-split? 41

context-element 15

convert-marker 26

convert-parent-wire-to-link 15

defined? 41

defined-indv-node? 41

derived? 41

disambiguate 22

do-elements 14

do-marked 25

downscan 49

end-kb-logging 12

english 21

eq-link? 41

eq-scan 49

free-marker 26

free-markers 26

function-node? 41

generic-indv-node? 41

get-cancel-marker 24

get-element-property 16

get-english-names 22

get-marker 26

get-set-node 32

get-type-node 32

has-link? 41

has-no-link? 41

incoming-a-elements 15

incoming-b-elements 15

incoming-c-elements 15

incoming-context-elements 15

incoming-parent-elements 15

incoming-split-elements 15

in-context 27

indv-map-node? 41

indv-node? 41

indv-role-node? 41

in-namespace 19

integer-node? 41

is-a-link? 41

is-not-a-link? 41

is-x-a-y? 42

kb-log 12

killed? 41

legal-marker? 24

legal-marker-pair? 24

link? 41

list-all-x-inverse-of-y 46

list-all-x-of-y 45

list-children 43

list-dictionary-entries 22

list-elements 43

list-equals 44

list-inferiors 44

list-instances 43

list-intersection 44

list-lowermost 44

list-marked 43

list-most-specific 45

list-namespaces 19

list-not-marked 43

list-parents 43

list-proper 44

list-rel 46

list-rel-inverse 47

list-roles 46

list-subtypes 43

list-superiors 44

list-supertypes 43

list-synonyms 22

list-uppermost 44

load-kb 8

lookup-definitions 21

lookup-element 19

lookup-element-or-defer 19

lookup-element-test 19

lowermost? 49

map-node? 41

mark 24

mark-all-x-inverse-of-y 46

mark-all-x-of-y 45

mark-boolean 25

mark-children 43

mark-equals 44

marker-bit 24

marker-count 24

marker-off? 25

marker-on? 25

mark-inferiors 44

mark-instances 43

mark-intersection 44

mark-lowermost 44

mark-most-specific 45

mark-named-elements 22

mark-parents 43

mark-proper 44

mark-rel 46

mark-rel-inverse 47

mark-role 45

mark-role-inverse 45

mark-roles 46

mark-subtypes 43

mark-superiors 44

mark-supertypes 43

mark-the-x-inverse-of-y 45

mark-the-x-of-y 45

mark-uppermost 44

new-cancel 35

new-complete-members 38

new-complete-split 37

new-complete-split-subtypes 38

new-context 27

new-defined-type 32

new-eq 34

new-float 31

new-function 31

new-indv 30

new-indv-role 38

new-integer 31

new-intersection-type 33

new-is-a 33

new-is-not-a 34

new-map 33

new-measure 52

new-members 38

new-not-statement 36

new-number 31

new-quantity 52

new-ratio 31

new-relation 36

new-split 37

new-split-subtypes 38

new-statement 36

new-string 31

new-struct 31

new-type 32

new-type-role 39

new-union-type 33

next-element 15

n-marker-pairs 24

n-markers 24

node? 41

node-value 31

not-eq-link? 41

not-statement? 41

number-node? 41

parent-element 15

previous-element 15

primitive-node? 41

process-deferred-connections 10

proper-indv-node? 41

push-element-property 16

ratio-node? 41

refresh-context 27

relation? 41

remove-context 40

remove-currently-loading-file 40

remove-element 40

remove-elements-after 40

remove-last-element 40

remove-last-loaded-file 40

role-node? 41

set-element-property 16

show-all-x-inverse-of-y 46

show-all-x-of-y 45

show-children 43

show-element-counts 47

show-elements 43

show-equals 44

show-inferiors 44

show-instances 43

show-intersection 44

show-lowermost 44

show-marked 43

show-marker-counts 47

show-most-specific 45

show-not-marked 43

show-ontology 47

show-parents 43

show-proper 44

show-rel 46

show-rel-inverse 47

show-roles 46

show-subtypes 43

show-superiors 44

show-supertypes 43

show-synonyms 22

show-uppermost 44

simple-is-x-a-y? 42

simple-is-x-eq-y? 42

split? 41

split-elements 15

splits-ok? 49

start-kb-logging 11

statement? 41

statement-true? 46

string-node? 41

struct-node? 41

symmetric? 41

the-x-inverse-of-y 46

the-x-of-y 45

the-x-of-y-is-a-z 39

the-x-of-y-is-not-a-z 39

the-x-role-of-y 46

transitive? 42

type-map-node? 41

type-node? 41

type-role-node? 41

unmark 24

unmark-boolean 25

uppermost? 49

upscan 48

which-split-member? 49

with-markers 26

x-is-a-y-of-z 39

x-is-not-a-y-of-z 39

x-is-not-the-y-of-z 39

x-is-the-y-of-z 39

\end{document}
